\documentclass[12pt]{article}
\usepackage[ukrainian]{babel}
\usepackage{fontspec}
% --- НАЛАШТУВАННЯ ШРИФТІВ ---
\setmainfont{Schoolbook-Regular.otf}[
    Path = ../../,
    BoldFont = Schoolbook-Bold.otf,
    ItalicFont = Schoolbook-Italic.otf,
    BoldItalicFont = Schoolbook-BoldItalic.otf
]
\usepackage[italic]{mathastext}
\usepackage[a4paper,margin=1.8cm,bottom=2cm,top=2cm]{geometry}
\usepackage{amsmath,amssymb}
\usepackage{tikz}
\usepackage{pgfplots}
\pgfplotsset{compat=1.16}
\usetikzlibrary{calc,patterns,angles,quotes}
\usepackage{xcolor,array,fancyhdr,enumitem,multicol}
\usepackage{colortbl}
\usepackage{tcolorbox}
\tcbuselibrary{skins}
\usepackage[hidelinks]{hyperref}
\usepackage[firstpage=false, text=@pvtr2525, scale=0.6, color=gray!12]{draftwatermark}
\definecolor{mainGreen}{RGB}{34, 120, 64}
\definecolor{yearOrange}{RGB}{220, 100, 30}
\definecolor{yearcolor}{RGB}{220, 100, 30}
\definecolor{titleBg}{RGB}{225, 240, 220}
\definecolor{instrBg}{RGB}{255, 240, 220}
\definecolor{instrBorder}{RGB}{220, 150, 80}
\pagestyle{fancy}
\fancyhf{}
\renewcommand{\headrulewidth}{0.4pt}
\renewcommand{\footrulewidth}{0.4pt}
\fancyhead[C]{\small\color{gray!80} Збірник НМТ \textendash{} сесії 23.05--5.06.2026}
\fancyhead[R]{\small\color{mainGreen}\textbf{\href{https://t.me/pvtr2525}{@pvtr2525}}}
\fancyfoot[L]{\small\color{gray!80}\href{https://t.me/pvtr2525}{ТГ: @pvtr2525}}
\fancyfoot[C]{\small\color{gray!80}\thepage}
\fancyfoot[R]{\small\color{gray!80}НМТ 2026}
% --- ТАБЛИЦІ ВІДПОВІДЕЙ ---
\newcommand{\answerTable}[5]{%
\par\vspace{0.15cm}
\noindent
\begin{tabular}{|*{5}{>{\centering\arraybackslash}m{3cm}|}}
\hline
\rule[-0.15cm]{0pt}{0.6cm}\textbf{А} & \rule[-0.15cm]{0pt}{0.6cm}\textbf{Б} & \rule[-0.15cm]{0pt}{0.6cm}\textbf{В} & \rule[-0.15cm]{0pt}{0.6cm}\textbf{Г} & \rule[-0.15cm]{0pt}{0.6cm}\textbf{Д} \\
\hline
\rule[-0.2cm]{0pt}{0.75cm}#1 & \rule[-0.2cm]{0pt}{0.75cm}#2 & \rule[-0.2cm]{0pt}{0.75cm}#3 & \rule[-0.2cm]{0pt}{0.75cm}#4 & \rule[-0.2cm]{0pt}{0.75cm}#5 \\
\hline
\end{tabular}
\par\vspace{0.25cm}
}
\newcommand{\answerTableTall}[5]{%
\par\vspace{0.15cm}
\noindent
\begin{tabular}{|*{5}{>{\centering\arraybackslash}m{3cm}|}}
\hline
\rule[-0.2cm]{0pt}{0.7cm}\textbf{А} & \rule[-0.2cm]{0pt}{0.7cm}\textbf{Б} & \rule[-0.2cm]{0pt}{0.7cm}\textbf{В} & \rule[-0.2cm]{0pt}{0.7cm}\textbf{Г} & \rule[-0.2cm]{0pt}{0.7cm}\textbf{Д} \\
\hline
\rule[-0.45cm]{0pt}{1.2cm}#1 & \rule[-0.45cm]{0pt}{1.2cm}#2 & \rule[-0.45cm]{0pt}{1.2cm}#3 & \rule[-0.45cm]{0pt}{1.2cm}#4 & \rule[-0.45cm]{0pt}{1.2cm}#5 \\
\hline
\end{tabular}
\par\vspace{0.25cm}
}
\newcommand{\instructionBox}[1]{%
\par\vspace{0.3cm}
\begin{tcolorbox}[colback=instrBg, colframe=instrBorder, boxrule=0.6pt, arc=2pt,
    left=8pt, right=8pt, top=4pt, bottom=4pt]
\centering\small\bfseries #1
\end{tcolorbox}
\par\vspace{0.15cm}
}
\newcommand{\sectionTitle}[1]{%
\par\vspace{0.4cm}
\begin{tcolorbox}[colback=titleBg, colframe=titleBg, boxrule=0pt, arc=8pt,
    halign=center, fontupper=\Large\bfseries\color{mainGreen}]
#1
\end{tcolorbox}
\par\vspace{0.3cm}
}
\newcommand{\task}[2]{\noindent\textbf{#1.}\ \ #2\par\vspace{0.2cm}}
% --- НМТ-СТИЛЬ ПОЛЕ ВІДПОВІДІ ---
\newcommand{\nmtAnswerBox}{%
\par\vspace{0.2cm}\noindent
Відповідь:\ \
\begin{tabular}{@{}*{4}{|>{\centering\arraybackslash}p{0.8cm}}|@{\hspace{0.1cm}\textbf{,}\hspace{0.1cm}}*{3}{|>{\centering\arraybackslash}p{0.8cm}}|@{}}
\hline
\rule[-0.2cm]{0pt}{0.7cm} & & & & & & \\
\hline
\end{tabular}
\par\vspace{0.3cm}
}
% --- СІТКА ДЛЯ МАТЧИНГУ ---
\newcommand{\matchingGrid}{%
    \begingroup
    \renewcommand{\arraystretch}{1.15}
    \setlength{\tabcolsep}{4pt}
    \footnotesize
    \begin{tabular}{r|c|c|c|c|c|}
         \multicolumn{1}{c}{} & \multicolumn{1}{c}{\textbf{А}} & \multicolumn{1}{c}{\textbf{Б}} & \multicolumn{1}{c}{\textbf{В}} & \multicolumn{1}{c}{\textbf{Г}} & \multicolumn{1}{c}{\textbf{Д}} \\ \cline{2-6}
         \textbf{1} & \phantom{X} & \phantom{X} & \phantom{X} & \phantom{X} & \phantom{X} \\ \cline{2-6}
         \textbf{2} & & & & & \\ \cline{2-6}
         \textbf{3} & & & & & \\ \cline{2-6}
    \end{tabular}
    \endgroup
}
% --- ДОДАТКОВІ МАКРОСИ ЗБІРНИКА ---
\newcommand{\taskBlock}[1]{%
\par\vspace{0.45cm}
\noindent{\color{mainGreen}\rule{\linewidth}{0.8pt}}\par\vspace{0.12cm}
\noindent{\small\bfseries\color{mainGreen}#1}\par\vspace{0.25cm}
}
\newcounter{zad}
\newcommand{\zadtask}[1]{\stepcounter{zad}\noindent\textbf{\thezad.}\ \ #1\par\vspace{0.2cm}}
\newcommand{\zadnum}{\stepcounter{zad}\textbf{\thezad.}\ \ }
\newcounter{ans}
\newcommand{\ansitem}[1]{\stepcounter{ans}\noindent\textbf{\theans.}\ #1\par\vspace{0.07cm}}
\newcommand{\nmtyear}[1]{\hfill{\small\color{yearOrange}(НМТ #1)}}
% --- ДОДАТКОВІ МАКРОСИ (розв'язки, зміст) ---
\tcbuselibrary{breakable}
\newcommand{\solution}[1]{%
\par\vspace{0.12cm}
\begin{tcolorbox}[colback=titleBg!45, colframe=mainGreen!55, boxrule=0.4pt, arc=2pt,
    left=6pt, right=6pt, top=3pt, bottom=3pt, breakable]
\footnotesize\textbf{Короткий розв'язок (оригінал).}\ #1
\end{tcolorbox}
\par\vspace{0.12cm}
}
\setcounter{tocdepth}{2}
\newcommand{\chapterTitle}[1]{%
\clearpage\phantomsection\addcontentsline{toc}{section}{#1}%
\sectionTitle{#1}}
\newcommand{\typeTitle}[1]{%
\phantomsection\addcontentsline{toc}{subsection}{\quad #1}%
\par\vspace{0.3cm}
\noindent{\large\bfseries\color{yearOrange}#1}\par\vspace{0.08cm}
\noindent{\color{yearOrange}\rule{\linewidth}{0.4pt}}\par\vspace{0.2cm}}
\newcommand{\ansTheme}[1]{\par\vspace{0.25cm}\noindent{\bfseries\color{mainGreen}#1}\par\vspace{0.12cm}}
\newcommand{\ansType}[1]{\par\vspace{0.1cm}\noindent{\itshape\color{yearOrange}#1}\par\vspace{0.08cm}}
\newcommand{\subTypeTitle}[1]{%
\par\vspace{0.35cm}
\phantomsection\addcontentsline{toc}{subsubsection}{\qquad #1}%
\noindent{\normalsize\bfseries\color{mainGreen}#1}\par\vspace{0.05cm}
\noindent{\color{mainGreen!60}\rule{\linewidth}{0.3pt}}\par\vspace{0.18cm}}
\begin{document}
\catcode`\"=12\relax
\thispagestyle{empty}
\vspace*{1.2cm}
\begin{center}
{\Large\bfseries\color{mainGreen} ЗБІРНИК НМТ-2026 з математики}\\[0.4cm]
{\Huge\bfseries\color{mainGreen} РОЗДІЛ 4. ПЛАНІМЕТРІЯ}\\[0.6cm]
{\large Оригінали сесій 23.05--5.06.2026 з короткими розв'язками та авторськими аналогами}\\[1cm]
{\large\color{mainGreen}\textbf{\href{https://t.me/pvtr2525}{Телеграм: @pvtr2525}}}
\end{center}
\clearpage
\phantomsection\addcontentsline{toc}{section}{Зміст}
\sectionTitle{ЗМІСТ}
{\hypersetup{linkcolor=black}\renewcommand{\contentsname}{}\vspace{-1cm}\tableofcontents}
\chapterTitle{РОЗДІЛ 4. ПЛАНІМЕТРІЯ}
\typeTitle{Завдання з вибором однієї відповіді (А--Д)}
\par\vspace{0.3cm}\noindent\textit{Розділ містить $84$ авторських аналоги до сесій НМТ-2026 з планіметрії: $63$ задачі з вибором однієї відповіді А--Д та $21$ задача на встановлення відповідності, розподілені на $5$ підтем за основними сюжетами шкільної геометрії.}
\par\noindent\textit{Підтеми \textbf{«Кути, відрізки та обчислення в трикутниках»} і \textbf{«Чотирикутники: довжини й площі»} тренують базові обчислення: знаходження кутів за сумою кутів трикутника, властивостями ромба й рівнобедреного трикутника, побудову прямокутних трикутників із кутами $30^\circ$, $45^\circ$, $60^\circ$ для обчислення сторін, діагоналей, висот, радіусів вписаного та описаного кіл. У підтемі \textbf{«Площі складених фігур та кільця»} зібрано задачі, де площу треба знайти як суму/різницю площ простих фігур або через формули круга та кільця.}
\par\noindent\textit{У підтемі \textbf{«Твердження та властивості»} та в розділі задач на відповідність важливо чітко відрізняти твердження, що виконуються \emph{завжди} (для будь-якого трикутника, чотирикутника, кола), від тих, що виконуються лише за додаткових умов: наприклад, медіана ділить трикутник на дві рівні за площею, але не за периметром; центр описаного кола лежить усередині лише для гострокутного трикутника; діагоналі ромба перпендикулярні, але \emph{не} рівні. Якщо в задачі сказано «довільний трикутник», то відповіді про рівнобедреність або перпендикулярність майже завжди хибні.}
\par\noindent\textit{Типові пастки: плутанина між кутом і його доповненням, забування про симетрію в задачах із кільцем чи басейном (чотири чверті $=$ одне коло), помилки в теоремі Піфагора у комбінованих фігурах, неточне розташування центра вписаного/описаного кола, ігнорування того, що бісектриса в \emph{рівнобедреному} трикутнику до основи є водночас висотою й медіаною, а в довільному --- ні.}
\par\noindent\textit{Під кожною оригінальною задачею (перша в трійці) розміщено кнопку \fbox{\footnotesize розв\textquotesingle язок $\rightarrow$}, що відкриває коментар у кінці розділу. Аналоги (друга й третя задачі) повторюють конфігурацію оригіналу зі зміненими параметрами або питанням --- їх корисно розв\textquotesingle язувати самостійно для закріплення методу.}
\par\vspace{0.3cm}
\typeTitle{Завдання з вибором однієї відповіді (А--Д)}
\subTypeTitle{Кути, відрізки та обчислення в трикутниках}
\taskBlock{Кути трикутника --- сесія 23.05, завдання №2}
\noindent
\begin{minipage}[c]{0.60\textwidth}
\zadnum Дано трикутник $ABC$ з кутами $\angle A=30^\circ$ та $\angle B=35^\circ$. Знайдіть градусну міру кута $C$. \nmtyear{2026}
\end{minipage}%
\hfill
\begin{minipage}[c]{0.38\textwidth}
\centering
\begin{tikzpicture}[scale=0.8, thick]
    \coordinate (A) at (0,0);
    \coordinate (C) at (4.5,0);
    \coordinate (B) at (1,2.8);
    \draw (A) -- (B) -- (C) -- cycle;
    \pic [draw, angle radius=0.5cm, "$30^\circ$"{right, shift={(0.15,0.15)}}] {angle = C--A--B};
    \pic [draw, angle radius=0.4cm, "$35^\circ$"{below, shift={(0.2,-0.1)}}] {angle = A--B--C};
    \pic [draw, angle radius=0.5cm, red, "?"{left, shift={(-0.15,0.15)}}] {angle = B--C--A};
    \node[below left] at (A) {$A$};
    \node[above] at (B) {$B$};
    \node[below right] at (C) {$C$};
\end{tikzpicture}
\end{minipage}
\par\vspace{0.2cm}
\answerTable{$75^\circ$}{$100^\circ$}{$115^\circ$}{$120^\circ$}{$30^\circ$}
\hypertarget{task:1}{}\par\vspace{0.05cm}\noindent\hyperlink{sol:1}{\fbox{\footnotesize розв'язок $\rightarrow$}}\par\vspace{0.2cm}
\vspace{0.3cm}
\noindent
\begin{minipage}[c]{0.60\textwidth}
\zadnum Дано трикутник $MNK$ з кутами $\angle M=40^\circ$ та $\angle N=65^\circ$. Знайдіть градусну міру кута $K$.
\end{minipage}%
\hfill
\begin{minipage}[c]{0.38\textwidth}
\centering
\begin{tikzpicture}[scale=0.8, thick]
    \coordinate (A) at (0,0);
    \coordinate (C) at (4.5,0);
    \coordinate (B) at (1.4,2.6);
    \draw (A) -- (B) -- (C) -- cycle;
    \pic [draw, angle radius=0.5cm, "$40^\circ$"{right, shift={(0.15,0.15)}}] {angle = C--A--B};
    \pic [draw, angle radius=0.4cm, "$65^\circ$"{below, shift={(0.2,-0.1)}}] {angle = A--B--C};
    \pic [draw, angle radius=0.5cm, red, "?"{left, shift={(-0.15,0.15)}}] {angle = B--C--A};
    \node[below left] at (A) {$M$};
    \node[above] at (B) {$N$};
    \node[below right] at (C) {$K$};
\end{tikzpicture}
\end{minipage}
\par\vspace{0.2cm}
\answerTable{$65^\circ$}{$75^\circ$}{$85^\circ$}{$105^\circ$}{$115^\circ$}
\vspace{0.3cm}
\noindent
\begin{minipage}[c]{0.60\textwidth}
\zadnum Дано трикутник $PQR$ з кутами $\angle P=25^\circ$ та $\angle Q=48^\circ$. Знайдіть градусну міру зовнішнього кута трикутника при вершині $R$.
\end{minipage}%
\hfill
\begin{minipage}[c]{0.38\textwidth}
\centering
\begin{tikzpicture}[scale=0.8, thick]
    \coordinate (A) at (0,0);
    \coordinate (C) at (4.5,0);
    \coordinate (B) at (1.2,2.4);
    \coordinate (F) at (6.0,0);
    \draw (A) -- (B) -- (C) -- cycle;
    \draw (C) -- (F);
    \pic [draw, angle radius=0.5cm, "$25^\circ$"{right, shift={(0.15,0.15)}}] {angle = C--A--B};
    \pic [draw, angle radius=0.4cm, "$48^\circ$"{below, shift={(0.2,-0.1)}}] {angle = A--B--C};
    \pic [draw, angle radius=0.5cm, red, "?"{above right, shift={(0.05,0.05)}}] {angle = F--C--B};
    \node[below left] at (A) {$P$};
    \node[above] at (B) {$Q$};
    \node[below] at (C) {$R$};
\end{tikzpicture}
\end{minipage}
\par\vspace{0.2cm}
\answerTable{$97^\circ$}{$103^\circ$}{$107^\circ$}{$117^\circ$}{$73^\circ$}
\vspace{0.3cm}
\newpage
\taskBlock{Планіметрія: кути ромба --- сесія 30.05, завдання №3}
\noindent
\begin{minipage}[c]{0.60\textwidth}
\zadnum Тупий кут ромба дорівнює $140^\circ$. Знайдіть кут, який утворює \textit{більша} діагональ цього ромба з його стороною.
\end{minipage}%
\hfill
\begin{minipage}[c]{0.38\textwidth}
\centering
\begin{tikzpicture}[scale=0.7, thick]
    \coordinate (A) at (-2.5, 0);
    \coordinate (B) at (0, 1.2);
    \coordinate (C) at (2.5, 0);
    \coordinate (D) at (0, -1.2);
    \draw[thick] (A) -- (B) -- (C) -- (D) -- cycle;
    \draw[thick] (A) -- (C);
    \node[left] at (A) {$A$};
    \node[above] at (B) {$B$};
    \node[right] at (C) {$C$};
    \node[below] at (D) {$D$};
\end{tikzpicture}
\end{minipage}
\par\vspace{0.2cm}
\answerTable{$20^\circ$}{$40^\circ$}{$50^\circ$}{$70^\circ$}{$110^\circ$}
\hypertarget{task:4}{}\par\vspace{0.05cm}\noindent\hyperlink{sol:4}{\fbox{\footnotesize розв'язок $\rightarrow$}}\par\vspace{0.2cm}
\noindent
\begin{minipage}[c]{0.60\textwidth}
\zadnum Тупий кут ромба дорівнює $110^\circ$. Знайдіть кут, який утворює \textit{більша} діагональ цього ромба з його стороною.
\end{minipage}%
\hfill
\begin{minipage}[c]{0.38\textwidth}
\centering
\begin{tikzpicture}[scale=0.7, thick]
    \coordinate (A) at (-2.2, 0);
    \coordinate (B) at (0, 1.5);
    \coordinate (C) at (2.2, 0);
    \coordinate (D) at (0, -1.5);
    \draw[thick] (A) -- (B) -- (C) -- (D) -- cycle;
    \draw[thick] (A) -- (C);
    \node[left] at (A) {$A$};
    \node[above] at (B) {$B$};
    \node[right] at (C) {$C$};
    \node[below] at (D) {$D$};
\end{tikzpicture}
\end{minipage}
\par\vspace{0.2cm}
\answerTable{$25^\circ$}{$35^\circ$}{$55^\circ$}{$70^\circ$}{$110^\circ$}
\noindent
\begin{minipage}[c]{0.60\textwidth}
\zadnum Тупий кут ромба дорівнює $150^\circ$. Знайдіть кут, який утворює \textit{менша} діагональ цього ромба з його стороною.
\end{minipage}%
\hfill
\begin{minipage}[c]{0.38\textwidth}
\centering
\begin{tikzpicture}[scale=0.7, thick]
    \coordinate (A) at (-2.7, 0);
    \coordinate (B) at (0, 1.0);
    \coordinate (C) at (2.7, 0);
    \coordinate (D) at (0, -1.0);
    \draw[thick] (A) -- (B) -- (C) -- (D) -- cycle;
    \draw[thick] (A) -- (C);
    \node[left] at (A) {$A$};
    \node[above] at (B) {$B$};
    \node[right] at (C) {$C$};
    \node[below] at (D) {$D$};
\end{tikzpicture}
\end{minipage}
\par\vspace{0.2cm}
\answerTable{$15^\circ$}{$30^\circ$}{$60^\circ$}{$75^\circ$}{$105^\circ$}
\newpage
\taskBlock{Півколо, поділене на рівні частини --- сесія 1.06, завдання №3}
\noindent
\begin{minipage}[c]{0.58\textwidth}
\zadnum Півколо розділене на $5$ рівних частин (див. рисунок). Визначте градусну міру кута $\alpha$.
\end{minipage}%
\hfill
\begin{minipage}[c]{0.40\textwidth}
\centering
\begin{tikzpicture}[scale=1.0, thick]
    \def\R{2}
    \coordinate (O) at (0,0);
    \draw[mainGreen!80!black] (\R,0) arc (0:180:\R);
    \draw[mainGreen!80!black] (-\R,0) -- (\R,0);
    \foreach \a in {36,72,108,144} {
        \draw[mainGreen!80!black] (0,0) -- (\a:\R);
    }
    \coordinate (P1) at (36:\R);
    \coordinate (P2) at (72:\R);
    \pic[draw=yearOrange, line width=0.8pt, "$\alpha$", angle radius=0.85cm,
         angle eccentricity=0.65, font=\small] {angle = P1--O--P2};
\end{tikzpicture}
\end{minipage}
\par\vspace{0.2cm}
\answerTable{$36^\circ$}{$32^\circ$}{$24^\circ$}{$42^\circ$}{$48^\circ$}
\hypertarget{task:7}{}\par\vspace{0.05cm}\noindent\hyperlink{sol:7}{\fbox{\footnotesize розв'язок $\rightarrow$}}\par\vspace{0.2cm}
\noindent
\begin{minipage}[c]{0.58\textwidth}
\zadnum Півколо розділене на $6$ рівних частин (див. рисунок). Визначте градусну міру кута $\alpha$.
\end{minipage}%
\hfill
\begin{minipage}[c]{0.40\textwidth}
\centering
\begin{tikzpicture}[scale=1.0, thick]
    \def\R{2}
    \coordinate (O) at (0,0);
    \draw[mainGreen!80!black] (\R,0) arc (0:180:\R);
    \draw[mainGreen!80!black] (-\R,0) -- (\R,0);
    \foreach \a in {30,60,90,120,150} {
        \draw[mainGreen!80!black] (0,0) -- (\a:\R);
    }
    \coordinate (P1) at (30:\R);
    \coordinate (P2) at (60:\R);
    \pic[draw=yearOrange, line width=0.8pt, "$\alpha$", angle radius=0.85cm,
         angle eccentricity=0.65, font=\small] {angle = P1--O--P2};
\end{tikzpicture}
\end{minipage}
\par\vspace{0.2cm}
\answerTable{$20^\circ$}{$25^\circ$}{$30^\circ$}{$36^\circ$}{$60^\circ$}
\noindent
\begin{minipage}[c]{0.58\textwidth}
\zadnum Півколо розділене на $6$ рівних частин (див. рисунок). Визначте градусну міру кута $\alpha$, що охоплює \textit{дві} сусідні частини.
\end{minipage}%
\hfill
\begin{minipage}[c]{0.40\textwidth}
\centering
\begin{tikzpicture}[scale=1.0, thick]
    \def\R{2}
    \coordinate (O) at (0,0);
    \draw[mainGreen!80!black] (\R,0) arc (0:180:\R);
    \draw[mainGreen!80!black] (-\R,0) -- (\R,0);
    \foreach \a in {30,60,90,120,150} {
        \draw[mainGreen!80!black] (0,0) -- (\a:\R);
    }
    \coordinate (P1) at (30:\R);
    \coordinate (P2) at (90:\R);
    \pic[draw=yearOrange, line width=0.8pt, "$\alpha$", angle radius=0.95cm,
         angle eccentricity=0.7, font=\small] {angle = P1--O--P2};
\end{tikzpicture}
\end{minipage}
\par\vspace{0.2cm}
\answerTable{$30^\circ$}{$45^\circ$}{$60^\circ$}{$72^\circ$}{$90^\circ$}
\newpage
\taskBlock{Кути рівнобедреного трикутника --- сесія 2.06, завдання №4}
\noindent
\begin{minipage}[c]{0.56\textwidth}
\zadnum На рисунку зображено клумбу у формі трикутника $ABC$, у якого $AB = BC$. Знайдіть градусну міру кута $B$, якщо він \textit{удвічі} більший за градусну міру кута $C$.
\end{minipage}%
\hfill
\begin{minipage}[c]{0.40\textwidth}
\centering
\begin{tikzpicture}[scale=0.5, thick]
    \coordinate (A) at (0,0);
    \coordinate (C) at (8,0);
    \coordinate (B) at (4,4);
    \draw[fill=mainGreen!15, draw=mainGreen!80!black, thick] (A) -- (B) -- (C) -- cycle;
    \node[below left] at (A) {$A$};
    \node[above] at (B) {$B$};
    \node[below right] at (C) {$C$};
    \pic[draw=yearOrange, line width=0.7pt, angle radius=0.3cm] {angle = C--A--B};
    \pic[draw=yearOrange, line width=0.7pt, angle radius=0.3cm] {angle = B--C--A};
    \pic[draw=yearOrange, line width=0.7pt, angle radius=0.3cm] {angle = A--B--C};
\end{tikzpicture}
\end{minipage}
\par\vspace{0.2cm}
\answerTable{$60^\circ$}{$30^\circ$}{$45^\circ$}{$90^\circ$}{$120^\circ$}
\hypertarget{task:10}{}\par\vspace{0.05cm}\noindent\hyperlink{sol:10}{\fbox{\footnotesize розв'язок $\rightarrow$}}\par\vspace{0.2cm}
\noindent
\begin{minipage}[c]{0.56\textwidth}
\zadnum На рисунку зображено ділянку у формі трикутника $ABC$, у якого $AB = BC$. Знайдіть градусну міру кута $B$, якщо він \textit{утричі} більший за градусну міру кута $C$.
\end{minipage}%
\hfill
\begin{minipage}[c]{0.40\textwidth}
\centering
\begin{tikzpicture}[scale=0.5, thick]
    \coordinate (A) at (0,0);
    \coordinate (C) at (8,0);
    \coordinate (B) at (4,2.9);
    \draw[fill=mainGreen!15, draw=mainGreen!80!black, thick] (A) -- (B) -- (C) -- cycle;
    \node[below left] at (A) {$A$};
    \node[above] at (B) {$B$};
    \node[below right] at (C) {$C$};
    \pic[draw=yearOrange, line width=0.7pt, angle radius=0.3cm] {angle = C--A--B};
    \pic[draw=yearOrange, line width=0.7pt, angle radius=0.3cm] {angle = B--C--A};
    \pic[draw=yearOrange, line width=0.7pt, angle radius=0.3cm] {angle = A--B--C};
\end{tikzpicture}
\end{minipage}
\par\vspace{0.2cm}
\answerTable{$36^\circ$}{$72^\circ$}{$90^\circ$}{$108^\circ$}{$120^\circ$}
\noindent
\begin{minipage}[c]{0.56\textwidth}
\zadnum На рисунку зображено газон у формі трикутника $ABC$, у якого $AB = BC$. Знайдіть градусну міру кута $A$ при основі, якщо кут $B$ \textit{на $30^\circ$ більший} за кут $A$.
\end{minipage}%
\hfill
\begin{minipage}[c]{0.40\textwidth}
\centering
\begin{tikzpicture}[scale=0.5, thick]
    \coordinate (A) at (0,0);
    \coordinate (C) at (8,0);
    \coordinate (B) at (4,4.77);
    \draw[fill=mainGreen!15, draw=mainGreen!80!black, thick] (A) -- (B) -- (C) -- cycle;
    \node[below left] at (A) {$A$};
    \node[above] at (B) {$B$};
    \node[below right] at (C) {$C$};
    \pic[draw=yearOrange, line width=0.7pt, angle radius=0.3cm] {angle = C--A--B};
    \pic[draw=yearOrange, line width=0.7pt, angle radius=0.3cm] {angle = B--C--A};
    \pic[draw=yearOrange, line width=0.7pt, angle radius=0.3cm] {angle = A--B--C};
\end{tikzpicture}
\end{minipage}
\par\vspace{0.2cm}
\answerTable{$50^\circ$}{$40^\circ$}{$65^\circ$}{$80^\circ$}{$100^\circ$}
\newpage
\taskBlock{Планіметрія: кути при перетині прямих --- сесія 3.06, завдання №4}
\noindent
\begin{minipage}[c]{0.58\textwidth}
\zadnum На рисунку схематично зображено підставку для планшета, що стоїть на горизонтальній поверхні (пряма $AD$). Точки $A$, $B$, $C$ і $K$ лежать на одній прямій. Визначте градусну міру кута $BAD$, якщо $\angle BKD = 128^\circ$, $\angle ADK = 68^\circ$.
\end{minipage}%
\hfill
\begin{minipage}[c]{0.40\textwidth}
\centering
\begin{tikzpicture}[scale=0.95, thick]
    \coordinate (A) at (0,0);
    \coordinate (D) at (2.4,0);
    \coordinate (B) at (1.7,3.0);
    \coordinate (K) at (1.25,2.2);
    \draw[thick] (A) -- (B);
    \draw[thick] (D) -- (K);
    \draw[thick] (A) -- (D);
    \node[above] at (B) {\scriptsize $B$};
    \node[below left] at (A) {\scriptsize $A$};
    \node[below right] at (D) {\scriptsize $D$};
    \node[left] at (K) {\scriptsize $K$};
    \pic[draw=yearOrange, line width=0.6pt, angle radius=0.4cm] {angle = D--K--B};
    \pic[draw=mainGreen, line width=0.6pt, angle radius=0.45cm] {angle = K--D--A};
    \pic[draw=yearOrange, line width=0.6pt, angle radius=0.4cm] {angle = D--A--B};
    \node[font=\tiny] at (2.2,2.05) {$128^\circ$};
    \node[font=\tiny] at (1.95,0.27) {$68^\circ$};
    \node[font=\tiny] at (0.52,0.18) {$?$};
\end{tikzpicture}
\end{minipage}
\par\vspace{0.2cm}
\answerTable{$66^\circ$}{$62^\circ$}{$52^\circ$}{$60^\circ$}{$50^\circ$}
\hypertarget{task:13}{}\par\vspace{0.05cm}\noindent\hyperlink{sol:13}{\fbox{\footnotesize розв'язок $\rightarrow$}}\par\vspace{0.2cm}
\noindent
\begin{minipage}[c]{0.58\textwidth}
\zadnum На рисунку схематично зображено підпірку для дзеркала, що стоїть на горизонтальній поверхні (пряма $AD$). Точки $A$, $B$, $C$ і $K$ лежать на одній прямій. Визначте градусну міру кута $BAD$, якщо $\angle BKD = 116^\circ$, $\angle ADK = 54^\circ$.
\end{minipage}%
\hfill
\begin{minipage}[c]{0.40\textwidth}
\centering
\begin{tikzpicture}[scale=0.95, thick]
    \coordinate (A) at (0,0);
    \coordinate (D) at (2.6,0);
    \coordinate (B) at (1.8,2.9);
    \coordinate (K) at (1.3,2.1);
    \draw[thick] (A) -- (B);
    \draw[thick] (D) -- (K);
    \draw[thick] (A) -- (D);
    \node[above] at (B) {\scriptsize $B$};
    \node[below left] at (A) {\scriptsize $A$};
    \node[below right] at (D) {\scriptsize $D$};
    \node[left] at (K) {\scriptsize $K$};
    \pic[draw=yearOrange, line width=0.6pt, angle radius=0.4cm] {angle = D--K--B};
    \pic[draw=mainGreen, line width=0.6pt, angle radius=0.45cm] {angle = K--D--A};
    \pic[draw=yearOrange, line width=0.6pt, angle radius=0.4cm] {angle = D--A--B};
    \node[font=\tiny] at (2.2,1.95) {$116^\circ$};
    \node[font=\tiny] at (2.1,0.27) {$54^\circ$};
    \node[font=\tiny] at (0.55,0.18) {$?$};
\end{tikzpicture}
\end{minipage}
\par\vspace{0.2cm}
\answerTable{$54^\circ$}{$58^\circ$}{$62^\circ$}{$64^\circ$}{$66^\circ$}
\noindent
\begin{minipage}[c]{0.58\textwidth}
\zadnum На рисунку схематично зображено підставку для книжки, що стоїть на горизонтальній поверхні (пряма $AD$). Точки $A$, $B$, $C$ і $K$ лежать на одній прямій. Визначте градусну міру кута $ADK$, якщо $\angle BKD = 140^\circ$, $\angle BAD = 85^\circ$.
\end{minipage}%
\hfill
\begin{minipage}[c]{0.40\textwidth}
\centering
\begin{tikzpicture}[scale=0.95, thick]
    \coordinate (A) at (0,0);
    \coordinate (D) at (2.5,0);
    \coordinate (B) at (1.6,3.0);
    \coordinate (K) at (1.2,2.25);
    \draw[thick] (A) -- (B);
    \draw[thick] (D) -- (K);
    \draw[thick] (A) -- (D);
    \node[above] at (B) {\scriptsize $B$};
    \node[below left] at (A) {\scriptsize $A$};
    \node[below right] at (D) {\scriptsize $D$};
    \node[left] at (K) {\scriptsize $K$};
    \pic[draw=yearOrange, line width=0.6pt, angle radius=0.4cm] {angle = D--K--B};
    \pic[draw=mainGreen, line width=0.6pt, angle radius=0.45cm] {angle = K--D--A};
    \pic[draw=yearOrange, line width=0.6pt, angle radius=0.4cm] {angle = D--A--B};
    \node[font=\tiny] at (2,2.05) {$140^\circ$};
    \node[font=\tiny] at (2.0,0.3) {$?$};
    \node[font=\tiny] at (0.52,0.18) {$85^\circ$};
\end{tikzpicture}
\end{minipage}
\par\vspace{0.2cm}
\answerTable{$55^\circ$}{$60^\circ$}{$65^\circ$}{$70^\circ$}{$75^\circ$}
\newpage
\taskBlock{Точка на відрізку --- сесія 4.06, завдання №3}
\zadtask{На відрізку $AB$ вибрано точку $K$ так, що відрізок $AK$ на $5$~\textit{см} більший за відрізок $BK$. Знайдіть довжину відрізка $AK$, якщо $AB = 21$~\textit{см}.}
\answerTable{$16$~\textit{см}}{$14$~\textit{см}}{$3$~\textit{см}}{$8$~\textit{см}}{$13$~\textit{см}}
\hypertarget{task:16}{}\par\vspace{0.05cm}\noindent\hyperlink{sol:16}{\fbox{\footnotesize розв'язок $\rightarrow$}}\par\vspace{0.2cm}
\zadtask{На відрізку $AB$ вибрано точку $K$ так, що відрізок $AK$ на $7$~\textit{см} більший за відрізок $BK$. Знайдіть довжину відрізка $AK$, якщо $AB = 25$~\textit{см}.}
\answerTable{$9$~\textit{см}}{$16$~\textit{см}}{$18$~\textit{см}}{$12$~\textit{см}}{$17$~\textit{см}}
\zadtask{На відрізку $MN$ вибрано точку $P$ так, що відрізок $MP$ на $4$~\textit{см} менший за відрізок $PN$. Знайдіть довжину відрізка $PN$, якщо $MN = 18$~\textit{см}.}
\answerTable{$7$~\textit{см}}{$14$~\textit{см}}{$9$~\textit{см}}{$11$~\textit{см}}{$13$~\textit{см}}
\newpage
\subTypeTitle{Чотирикутники: довжини й площі}
\taskBlock{Квадрат, радіус описаного кола --- сесія 1.06, завдання №10}
\noindent
\begin{minipage}[c]{0.60\textwidth}
\zadnum На рисунку зображено квадрат $ABCD$. На стороні $AD$ вибрано точку $K$ так, що $\angle BKD = 120^\circ$. Визначте радіус кола, описаного навколо цього квадрата, якщо $BK = 12$ см.
\end{minipage}%
\hfill
\begin{minipage}[c]{0.38\textwidth}
\centering
\begin{tikzpicture}[scale=0.85, thick]
    \coordinate (A) at (0,0);
    \coordinate (D) at (3,0);
    \coordinate (C) at (3,3);
    \coordinate (B) at (0,3);
    \coordinate (K) at (1.7,0);
    \draw[thick] (A) -- (B) -- (C) -- (D) -- cycle;
    \draw[thick] (B) -- (K);
    \fill (K) circle (1.8pt);
    \pic[draw=yearOrange, line width=0.8pt, "$120^\circ$", angle radius=0.3cm,
         angle eccentricity=1.85, font=\scriptsize] {angle = D--K--B};
    \node[below left] at (A) {$A$};
    \node[above left] at (B) {$B$};
    \node[above right] at (C) {$C$};
    \node[below right] at (D) {$D$};
    \node[below] at (K) {$K$};
\end{tikzpicture}
\end{minipage}
\par\vspace{0.2cm}
\answerTable{$6\sqrt{2}$ см}{$3\sqrt{3}$ см}{$6\sqrt{3}$ см}{$3\sqrt{2}$ см}{$3\sqrt{6}$ см}
\hypertarget{task:19}{}\par\vspace{0.05cm}\noindent\hyperlink{sol:19}{\fbox{\footnotesize розв'язок $\rightarrow$}}\par\vspace{0.2cm}
\noindent
\begin{minipage}[c]{0.60\textwidth}
\zadnum На рисунку зображено прямокутник $ABCD$. На стороні $AD$ вибрано точку $K$ так, що $\angle BKD = 135^\circ$. Визначте довжину меншої сторони цього прямокутника, якщо $BK = 8$~см.
\end{minipage}%
\hfill
\begin{minipage}[c]{0.38\textwidth}
\centering
\begin{tikzpicture}[scale=0.85, thick, line join=round]
    \coordinate (A) at (0,0);
    \coordinate (D) at (4,0);
    \coordinate (C) at (4,2.5);
    \coordinate (B) at (0,2.5);
    \coordinate (K) at (2.5,0);
    \draw (A) -- (B) -- (C) -- (D) -- cycle;
    \draw (B) -- (K);
    \fill (K) circle (2pt);
    \pic[draw=yearOrange, line width=0.8pt, "$135^\circ$", angle radius=0.45cm,
         angle eccentricity=1.6, font=\scriptsize] {angle = D--K--B};
    \node[below left] at (A) {$A$};
    \node[above left] at (B) {$B$};
    \node[above right] at (C) {$C$};
    \node[below right] at (D) {$D$};
    \node[below] at (K) {$K$};
\end{tikzpicture}
\end{minipage}
\par\vspace{0.2cm}
\answerTable{$2\sqrt{2}$~см}{$4$~см}{$4\sqrt{2}$~см}{$8$~см}{$8\sqrt{2}$~см}
\noindent
\begin{minipage}[c]{0.60\textwidth}
\zadnum На рисунку зображено квадрат $ABCD$ зі стороною $16$~см. Точка $K$~--- середина сторони $AD$. Визначте радіус кола, описаного навколо трикутника $BKC$.
\end{minipage}%
\hfill
\begin{minipage}[c]{0.38\textwidth}
\centering
\begin{tikzpicture}[scale=0.7, thick]
    \coordinate (A) at (0,0);
    \coordinate (D) at (3,0);
    \coordinate (C) at (3,3);
    \coordinate (B) at (0,3);
    \coordinate (K) at (1.5,0);
    \draw[thick] (A) -- (B) -- (C) -- (D) -- cycle;
    \draw[thick, mainGreen] (B) -- (K) -- (C) -- cycle;
    \fill (K) circle (1.8pt);
    \node[below left] at (A) {$A$};
    \node[above left] at (B) {$B$};
    \node[above right] at (C) {$C$};
    \node[below right] at (D) {$D$};
    \node[below] at (K) {$K$};
\end{tikzpicture}
\end{minipage}
\par\vspace{0.2cm}
\answerTable{$8$ см}{$10$ см}{$8\sqrt{2}$ см}{$4\sqrt{5}$ см}{$5\sqrt{5}$ см}
\newpage
\taskBlock{Площа рівнобічної трапеції --- сесія 2.06, завдання №3}
\noindent
\begin{minipage}[c]{0.54\textwidth}
\zadnum У рівнобічній трапеції $ABCD$ бічна сторона дорівнює $6$, а гострий кут при більшій основі дорівнює $60^\circ$. З вершин меншої основи $B$ і $C$ проведено відрізки до спільної точки $P$ на більшій основі так, що $BP \perp AB$ і $CP \perp CD$ (див. рисунок). Знайдіть площу трапеції.
\end{minipage}%
\hfill
\begin{minipage}[c]{0.44\textwidth}
\centering
\begin{tikzpicture}[scale=0.27, thick]
    \coordinate (A) at (0,0);
    \coordinate (D) at (24,0);
    \coordinate (B) at (3,5.196);
    \coordinate (C) at (21,5.196);
    \coordinate (P) at (12,0);
    \draw[thick] (A) -- (D) -- (C) -- (B) -- cycle;
    \draw[thick, mainGreen] (B) -- (P) -- (C);
    \pic[draw=yearOrange, line width=0.7pt, angle radius=0.3cm] {right angle = A--B--P};
    \pic[draw=yearOrange, line width=0.7pt, angle radius=0.3cm] {right angle = P--C--D};
    \pic[draw=yearOrange, line width=0.7pt, pic text={\scriptsize $60^\circ$}, angle radius=0.3cm, angle eccentricity=1.9] {angle = D--A--B};
    \node[below left] at (A) {$A$};
    \node[above left] at (B) {$B$};
    \node[above right] at (C) {$C$};
    \node[below right] at (D) {$D$};
    \node[below] at (P) {$P$};
    \node[above left, font=\scriptsize] at ($(B)!0.5!(A)$) {$6$};
\end{tikzpicture}
\end{minipage}
\par\vspace{0.2cm}
\answerTable{$63$}{$252$}{$126$}{$126\sqrt{3}$}{$63\sqrt{3}$}
\hypertarget{task:22}{}\par\vspace{0.05cm}\noindent\hyperlink{sol:22}{\fbox{\footnotesize розв'язок $\rightarrow$}}\par\vspace{0.2cm}
\noindent
\begin{minipage}[c]{0.54\textwidth}
\zadnum У рівнобічній трапеції $ABCD$ бічна сторона дорівнює $6$, а гострий кут при більшій основі дорівнює $30^\circ$. З вершин меншої основи $B$ і $C$ проведено відрізки до спільної точки $P$ на більшій основі так, що $BP \perp AB$ і $CP \perp CD$ (див. рисунок). Знайдіть площу трапеції.
\end{minipage}%
\hfill
\begin{minipage}[c]{0.44\textwidth}
\centering
\begin{tikzpicture}[scale=0.42, thick]
    \coordinate (A) at (0,0);
    \coordinate (D) at (13.86,0);
    \coordinate (B) at (5.196,3);
    \coordinate (C) at (8.66,3);
    \coordinate (P) at (6.93,0);
    \draw[thick] (A) -- (D) -- (C) -- (B) -- cycle;
    \draw[thick, mainGreen] (B) -- (P) -- (C);
    \pic[draw=yearOrange, line width=0.7pt, angle radius=0.3cm] {right angle = A--B--P};
    \pic[draw=yearOrange, line width=0.7pt, angle radius=0.3cm] {right angle = P--C--D};
    \pic[draw=yearOrange, line width=0.7pt, pic text={\scriptsize $30^\circ$}, angle radius=0.5cm, angle eccentricity=1.9] {angle = D--A--B};
    \node[below left] at (A) {$A$};
    \node[above left] at (B) {$B$};
    \node[above right] at (C) {$C$};
    \node[below right] at (D) {$D$};
    \node[below] at (P) {$P$};
    \node[above left, font=\scriptsize] at ($(B)!0.5!(A)$) {$6$};
\end{tikzpicture}
\end{minipage}
\par\vspace{0.2cm}
\answerTable{$30\sqrt{3}$}{$15\sqrt{3}$}{$15$}{$10\sqrt{3}$}{$45$}
\noindent
\begin{minipage}[c]{0.54\textwidth}
\zadnum У прямокутній трапеції $ABCD$ ($\angle A = \angle B = 90^\circ$) менша основа $BC = 5$, висота $AB = 4$, а гострий кут $\angle ADC = 45^\circ$ (див. рисунок). Знайдіть площу трапеції.
\end{minipage}%
\hfill
\begin{minipage}[c]{0.44\textwidth}
\centering
\begin{tikzpicture}[scale=0.45, thick]
    \coordinate (A) at (0,0);
    \coordinate (B) at (0,4);
    \coordinate (C) at (5,4);
    \coordinate (D) at (9,0);
    \draw[thick] (A) -- (B) -- (C) -- (D) -- cycle;
    \pic[draw=yearOrange, line width=0.7pt, angle radius=0.3cm] {right angle = D--A--B};
    \pic[draw=yearOrange, line width=0.7pt, angle radius=0.3cm] {right angle = A--B--C};
    \pic[draw=yearOrange, line width=0.7pt, pic text={\scriptsize $45^\circ$}, angle radius=0.45cm, angle eccentricity=1.8] {angle = C--D--A};
    \node[below left] at (A) {$A$};
    \node[above left] at (B) {$B$};
    \node[above right] at (C) {$C$};
    \node[below right] at (D) {$D$};
    \node[above, font=\scriptsize] at (2.5,4) {$5$};
    \node[left, font=\scriptsize] at (0,2) {$4$};
\end{tikzpicture}
\end{minipage}
\par\vspace{0.2cm}
\answerTable{$36$}{$28$}{$20$}{$14$}{$32$}
\newpage
\taskBlock{Планіметрія: прямокутник і вписане коло --- сесія 3.06, завдання №14}
\noindent
\begin{minipage}[c]{0.56\textwidth}
\zadnum На рисунку зображено прямокутник $ABCD$ та коло з центром $O$, яке дотикається до сторін $AB$, $AD$, а також сторони $BC$ у точці $K$. Визначте площу чотирикутника $OKCD$, якщо $OD = 24$ см, $\angle ADO = 30^\circ$.
\end{minipage}%
\hfill
\begin{minipage}[c]{0.40\textwidth}
\centering
\begin{tikzpicture}[scale=0.42, thick]
    \coordinate (A) at (0,0);
    \coordinate (B) at (0,5);
    \coordinate (Cc) at (8.7,5);
    \coordinate (D) at (8.7,0);
    \coordinate (O) at (2.5,2.5);
    \coordinate (K) at (2.5,5);
    \draw[thick] (A) -- (B) -- (Cc) -- (D) -- cycle;
    \draw[thick, mainGreen] (O) circle (2.5);
    \draw[thick] (O) -- (D);
    \draw[thick] (O) -- (K);
    \fill (O) circle (2.4pt); \node[left, font=\scriptsize] at (O) {$O$};
    \fill (K) circle (2.4pt); \node[above, font=\scriptsize] at (K) {$K$};
    \node[below left, font=\scriptsize] at (A) {$A$};
    \node[above left, font=\scriptsize] at (B) {$B$};
    \node[above right, font=\scriptsize] at (Cc) {$C$};
    \node[below right, font=\scriptsize] at (D) {$D$};
    \pic[draw, angle radius=0.9cm] {angle = O--D--A};
    \node[font=\tiny] at (5.9,0.5) {$30^\circ$};
\end{tikzpicture}
\end{minipage}
\par\vspace{0.2cm}
\answerTable{$288$ см$^2$}{$576$ см$^2$}{$216\sqrt{3}$ см$^2$}{$288\sqrt{3}$ см$^2$}{$576\sqrt{3}$ см$^2$}
\hypertarget{task:25}{}\par\vspace{0.05cm}\noindent\hyperlink{sol:25}{\fbox{\footnotesize розв'язок $\rightarrow$}}\par\vspace{0.2cm}
\noindent
\begin{minipage}[c]{0.56\textwidth}
\zadnum На рисунку зображено прямокутник $ABCD$ та коло з центром $O$, яке дотикається до сторін $AB$, $AD$, а також сторони $BC$ у точці $K$. Визначте площу чотирикутника $OKCD$, якщо $OD = 12\sqrt{2}$ см, $\angle ADO = 45^\circ$.
\end{minipage}%
\hfill
\begin{minipage}[c]{0.40\textwidth}
\centering
\begin{tikzpicture}[scale=0.42, thick]
    \coordinate (A) at (0,0);
    \coordinate (B) at (0,5);
    \coordinate (Cc) at (5.5,5);
    \coordinate (D) at (5.5,0);
    \coordinate (O) at (2.5,2.5);
    \coordinate (K) at (2.5,5);
    \draw[thick] (A) -- (B) -- (Cc) -- (D) -- cycle;
    \draw[thick, mainGreen] (O) circle (2.5);
    \draw[thick] (O) -- (D);
    \draw[thick] (O) -- (K);
    \fill (O) circle (2.4pt); \node[left, font=\scriptsize] at (O) {$O$};
    \fill (K) circle (2.4pt); \node[above, font=\scriptsize] at (K) {$K$};
    \node[below left, font=\scriptsize] at (A) {$A$};
    \node[above left, font=\scriptsize] at (B) {$B$};
    \node[above right, font=\scriptsize] at (Cc) {$C$};
    \node[below right, font=\scriptsize] at (D) {$D$};
    \pic[draw, angle radius=0.8cm] {angle = O--D--A};
    \node[font=\tiny] at (3.6,0.5) {$45^\circ$};
\end{tikzpicture}
\end{minipage}
\par\vspace{0.2cm}
\answerTable{$144\sqrt{2}$ см$^2$}{$108$ см$^2$}{$216$ см$^2$}{$216\sqrt{2}$ см$^2$}{$432$ см$^2$}
\noindent
\begin{minipage}[c]{0.56\textwidth}
\zadnum На рисунку зображено прямокутник $ABCD$, у якому $AB = 5$~см, $AD = 12$~см. Коло з центром $O$ вписане в трикутник $ABD$, тобто дотикається сторін $AB$, $AD$ та діагоналі $BD$. Знайдіть радіус цього кола.
\end{minipage}%
\hfill
\begin{minipage}[c]{0.40\textwidth}
\centering
\begin{tikzpicture}[scale=0.32, thick]
    \coordinate (A) at (0,0);
    \coordinate (B) at (0,5);
    \coordinate (Cc) at (12,5);
    \coordinate (D) at (12,0);
    \coordinate (O) at (2,2);
    \draw[thick] (A) -- (B) -- (Cc) -- (D) -- cycle;
    \draw[thick, mainGreen] (B) -- (D);
    \draw[thick, mainGreen!80!black] (O) circle (2);
    \fill (O) circle (3pt); \node[right=2pt, font=\scriptsize] at (O) {$O$};
    \node[below left, font=\scriptsize] at (A) {$A$};
    \node[above left, font=\scriptsize] at (B) {$B$};
    \node[above right, font=\scriptsize] at (Cc) {$C$};
    \node[below right, font=\scriptsize] at (D) {$D$};
\end{tikzpicture}
\end{minipage}
\par\vspace{0.2cm}
\answerTable{$2$ см}{$2{,}5$ см}{$\dfrac{30}{13}$ см}{$3$ см}{$\dfrac{13}{6}$ см}
\newpage
\taskBlock{Площа паралелограма --- сесія 4.06, завдання №14}
\noindent
\begin{minipage}[c]{0.54\textwidth}
\zadnum Точка $K$~--- середина висоти $BM$ паралелограма $ABCD$ (див. рисунок). $AB = 8\sqrt{3}$~см, $DK = 12$~см, $\angle ADK = 30^\circ$. Визначте площу паралелограма $ABCD$.
\end{minipage}%
\hfill
\begin{minipage}[c]{0.44\textwidth}
\centering
\begin{tikzpicture}[scale=0.55, thick]
    \coordinate (A) at (0,0);
    \coordinate (M) at (2.5,0);
    \coordinate (D) at (7.5,0);
    \coordinate (B) at (2.5,4);
    \coordinate (C) at (10,4);
    \coordinate (K) at (2.5,2);
    \draw[thick] (A) -- (B) -- (C) -- (D) -- cycle;
    \draw[thick] (B) -- (M);
    \draw[thick] (K) -- (D);
    \fill (K) circle (2.5pt);
    \draw (M) ++(0,0.3) -- ++(0.3,0) -- ++(0,-0.3);
    \pic[draw, angle radius=0.9cm] {angle = K--D--M};
    \node at ($(D)+(170:2.5)$) {\scriptsize $30^\circ$};
    \node[below left] at (A) {$A$};
    \node[above left] at (B) {$B$};
    \node[above right] at (C) {$C$};
    \node[below right] at (D) {$D$};
    \node[below] at (M) {$M$};
    \node[left] at (K) {$K$};
\end{tikzpicture}
\end{minipage}
\par\vspace{0.2cm}
\answerTable{$120\sqrt{3}$~см$^2$}{$72\sqrt{3}$~см$^2$}{$96\sqrt{3}$~см$^2$}{$120$~см$^2$}{$60\sqrt{3}$~см$^2$}
\hypertarget{task:28}{}\par\vspace{0.05cm}\noindent\hyperlink{sol:28}{\fbox{\footnotesize розв'язок $\rightarrow$}}\par\vspace{0.2cm}
\noindent
\begin{minipage}[c]{0.54\textwidth}
\zadnum Точка $K$~--- середина висоти $BM$ паралелограма $ABCD$ (див. рисунок). $AB = 12$~см, $DK = 6$~см, $\angle ADK = 60^\circ$. Визначте площу паралелограма $ABCD$.
\end{minipage}%
\hfill
\begin{minipage}[c]{0.44\textwidth}
\centering
\begin{tikzpicture}[scale=0.55, thick]
    \coordinate (A) at (0,0);
    \coordinate (M) at (3.2,0);
    \coordinate (D) at (5.2,0);
    \coordinate (B) at (3.2,4.4);
    \coordinate (C) at (8.4,4.4);
    \coordinate (K) at (3.2,2.2);
    \draw[thick] (A) -- (B) -- (C) -- (D) -- cycle;
    \draw[thick] (B) -- (M);
    \draw[thick] (K) -- (D);
    \fill (K) circle (2.5pt);
    \draw (M) ++(0,0.3) -- ++(0.3,0) -- ++(0,-0.3);
    \pic[draw, angle radius=0.55cm] {angle = K--D--M};
    \node at ($(D)+(150:1.4)$) {\scriptsize $60^\circ$};
    \node[below left] at (A) {$A$};
    \node[above left] at (B) {$B$};
    \node[above right] at (C) {$C$};
    \node[below right] at (D) {$D$};
    \node[below] at (M) {$M$};
    \node[left] at (K) {$K$};
\end{tikzpicture}
\end{minipage}
\par\vspace{0.2cm}
\answerTable{$27\sqrt{3}$~см$^2$}{$36\sqrt{3}$~см$^2$}{$54$~см$^2$}{$54\sqrt{3}$~см$^2$}{$108\sqrt{3}$~см$^2$}
\noindent
\begin{minipage}[c]{0.54\textwidth}
\zadnum У рівнобедреному трикутнику $ABC$ ($AB = BC$) проведено висоту $BH$ з вершини $B$ на основу $AC$. Точка $K$~--- середина висоти $BH$ (див. рисунок). $AC = 12$~см, $\angle ACK = 30^\circ$. Визначте площу трикутника $ABC$.
\end{minipage}%
\hfill
\begin{minipage}[c]{0.44\textwidth}
\centering
\begin{tikzpicture}[scale=0.55, thick]
    \coordinate (A) at (-3,0);
    \coordinate (C) at (3,0);
    \coordinate (H) at (0,0);
    \coordinate (B) at (0,4);
    \coordinate (K) at (0,2);
    \draw[thick] (A) -- (B) -- (C) -- cycle;
    \draw[thick] (B) -- (H);
    \draw[thick, mainGreen] (C) -- (K);
    \fill (K) circle (2.5pt);
    \draw (H) ++(0,0.3) -- ++(0.3,0) -- ++(0,-0.3);
    \pic[draw=yearOrange, line width=0.7pt, pic text={\scriptsize $30^\circ$}, angle radius=0.55cm, angle eccentricity=1.55] {angle = K--C--A};
    \node[below left] at (A) {$A$};
    \node[above] at (B) {$B$};
    \node[below right] at (C) {$C$};
    \node[below] at (H) {$H$};
    \node[left] at (K) {$K$};
\end{tikzpicture}
\end{minipage}
\par\vspace{0.2cm}
\answerTable{$12\sqrt{3}$~см$^2$}{$48\sqrt{3}$~см$^2$}{$24\sqrt{3}$~см$^2$}{$36\sqrt{3}$~см$^2$}{$24$~см$^2$}
\newpage
\taskBlock{Площа прямокутного трикутника у прямокутнику --- сесія 5.06, завдання №14}
\noindent
\begin{minipage}[c]{0.56\textwidth}
\zadnum У прямокутник $ABCD$ вписано прямокутний трикутник $AKM$ так, як зображено на рисунку. Обчисліть площу трикутника $AKM$, якщо $AK = 12$ см, $AB = 15$ см, $\angle KAD = 30^\circ$.
\end{minipage}%
\hfill
\begin{minipage}[c]{0.40\textwidth}
\centering
\begin{tikzpicture}[scale=0.32, thick]
    \coordinate (A) at (0,0);
    \coordinate (B) at (0,15);
    \coordinate (C) at (11,15);
    \coordinate (D) at (11,0);
    \coordinate (K) at (11,6);
    \coordinate (M) at (5.2,15);
    \draw[thick] (A) -- (B) -- (C) -- (D) -- cycle;
    \draw[thick] (A) -- (K) -- (M) -- cycle;
    \pic[draw=yearOrange, line width=0.6pt, angle radius=0.7cm] {right angle = A--K--M};
    \draw[thick] (K) -- (A) -- (D) -- cycle;
    \pic[draw=yearOrange, line width=0.6pt, angle radius=0.3cm] { angle = D--A--K};
    \node[below left, font=\scriptsize] at (A) {$A$};
    \node[above left, font=\scriptsize] at (B) {$B$};
    \node[above, font=\scriptsize] at (M) {$M$};
    \node[above right, font=\scriptsize] at (C) {$C$};
    \node[right, font=\scriptsize] at (K) {$K$};
    \node[below right, font=\scriptsize] at (D) {$D$};
    \node[font=\tiny] at (1.9,0.5) {$30^\circ$};
\end{tikzpicture}
\end{minipage}
\par\vspace{0.2cm}
\answerTable{$36$ см$^2$}{$96$ см$^2$}{$72$ см$^2$}{$72\sqrt{3}$ см$^2$}{$36\sqrt{3}$ см$^2$}
\hypertarget{task:31}{}\par\vspace{0.05cm}\noindent\hyperlink{sol:31}{\fbox{\footnotesize розв'язок $\rightarrow$}}\par\vspace{0.2cm}
\noindent
\begin{minipage}[c]{0.56\textwidth}
\zadnum У прямокутник $ABCD$ вписано прямокутний трикутник $AKM$ так, як зображено на рисунку. Знайдіть довжину катета $KM$, якщо $AK = 8$ см, $AB = 10$ см, $\angle KAD = 30^\circ$.
\end{minipage}%
\hfill
\begin{minipage}[c]{0.40\textwidth}
\centering
\begin{tikzpicture}[scale=0.45, thick]
    \coordinate (A) at (0,0);
    \coordinate (B) at (0,10);
    \coordinate (C) at (7.4,10);
    \coordinate (D) at (7.4,0);
    \coordinate (K) at (7.4,4);
    \coordinate (M) at (3.6,10);
    \draw[thick] (A) -- (B) -- (C) -- (D) -- cycle;
    \draw[thick] (A) -- (K) -- (M) -- cycle;
    \pic[draw=yearOrange, line width=0.6pt, angle radius=0.4cm] {right angle = A--K--M};
    \pic[draw=yearOrange, line width=0.6pt, angle radius=0.4cm] { angle = D--A--K};
    \node[below left, font=\scriptsize] at (A) {$A$};
    \node[above left, font=\scriptsize] at (B) {$B$};
    \node[above, font=\scriptsize] at (M) {$M$};
    \node[above right, font=\scriptsize] at (C) {$C$};
    \node[right, font=\scriptsize] at (K) {$K$};
    \node[below right, font=\scriptsize] at (D) {$D$};
    \node[font=\tiny] at (1.5,0.4) {$30^\circ$};
\end{tikzpicture}
\end{minipage}
\par\vspace{0.2cm}
\answerTable{$2\sqrt{3}$ см}{$4\sqrt{3}$ см}{$6$ см}{$4\sqrt{2}$ см}{$8\sqrt{3}$ см}
\noindent
\begin{minipage}[c]{0.56\textwidth}
\zadnum У прямокутний трикутник $ABC$ з прямим кутом при $A$ і катетами $AB = 6$~см та $AC = 12$~см вписано квадрат $APQR$ так, що дві його сторони лежать на катетах $AC$ і $AB$ (див. рисунок). Знайдіть площу квадрата $APQR$.
\end{minipage}%
\hfill
\begin{minipage}[c]{0.40\textwidth}
\centering
\begin{tikzpicture}[scale=0.32, thick]
    \coordinate (A) at (0,0);
    \coordinate (B) at (0,6);
    \coordinate (C) at (12,0);
    \coordinate (P) at (4,0);
    \coordinate (Q) at (4,4);
    \coordinate (R) at (0,4);
    \draw[thick] (A) -- (B) -- (C) -- cycle;
    \draw[thick, fill=mainGreen!15] (A) -- (P) -- (Q) -- (R) -- cycle;
    \pic[draw=yearOrange, line width=0.6pt, angle radius=0.45cm] {right angle = C--A--B};
    \node[below left, font=\scriptsize] at (A) {$A$};
    \node[above, font=\scriptsize] at (B) {$B$};
    \node[below right, font=\scriptsize] at (C) {$C$};
    \node[below, font=\scriptsize] at (P) {$P$};
    \node[above right, font=\scriptsize] at (Q) {$Q$};
    \node[left, font=\scriptsize] at (R) {$R$};
\end{tikzpicture}
\end{minipage}
\par\vspace{0.2cm}
\answerTable{$16$ см$^2$}{$9$ см$^2$}{$25$ см$^2$}{$36$ см$^2$}{$12$ см$^2$}
\newpage
\subTypeTitle{Площі складених фігур та кільця}
\taskBlock{Площа складеної фігури --- сесія 23.05, завдання №13}
\noindent
\begin{minipage}[c]{0.55\textwidth}
\zadnum Для виготовлення картонної моделі трафарету $ABCDEG$ Андрійко наклеїв квадрат зі стороною $6$~см на прямокутник зі сторонами $4$ і $9$~см (див. рисунок). Спільною частиною квадрата та прямокутника є рівнобедрений трикутник, дві вершини якого збігаються з вершинами прямокутника. Обчисліть площу трафарету $ABCDEG$. \nmtyear{2026}
\end{minipage}%
\hfill
\begin{minipage}[c]{0.43\textwidth}
\centering
\begin{tikzpicture}[scale=0.45, thick]
    \coordinate (A) at (-3,0);
    \coordinate (B) at (0,3);
    \coordinate (Rt) at (3,0);
    \coordinate (G) at (0,-3);
    \coordinate (C) at (1,2);
    \coordinate (K) at (1,-2);
    \coordinate (D) at (7,2);
    \coordinate (E) at (7,-2);
    \fill[gray!35] (C) -- (Rt) -- (K) -- cycle;
    \draw (A) -- (B) -- (Rt) -- (G) -- cycle;
    \draw (C) -- (D) -- (E) -- (K);
    \draw (C) -- (K);
    \node[left] at (A) {$A$};
    \node[above] at (B) {$B$};
    \node[above right] at (C) {$C$};
    \node[above right] at (D) {$D$};
    \node[below right] at (E) {$E$};
    \node[below] at (G) {$G$};
\end{tikzpicture}
\end{minipage}
\par\vspace{0.2cm}
\answerTable{$68$ см$^2$}{$64$ см$^2$}{$72$ см$^2$}{$70$ см$^2$}{$76$ см$^2$}
\hypertarget{task:34}{}\par\vspace{0.05cm}\noindent\hyperlink{sol:34}{\fbox{\footnotesize розв'язок $\rightarrow$}}\par\vspace{0.2cm}
\vspace{0.3cm}
\noindent
\begin{minipage}[c]{0.55\textwidth}
\zadnum Для виготовлення картонної моделі трафарету $ABCDEG$ Оля наклеїла квадрат зі стороною $10$~см на прямокутник зі сторонами $8$ і $12$~см (див. рисунок). Спільною частиною квадрата та прямокутника є рівнобедрений трикутник, дві вершини якого збігаються з вершинами прямокутника. Обчисліть площу трафарету $ABCDEG$.
\end{minipage}%
\hfill
\begin{minipage}[c]{0.43\textwidth}
\centering
\begin{tikzpicture}[scale=0.33, thick]
    \coordinate (A) at (-5,0);
    \coordinate (B) at (0,5);
    \coordinate (Rt) at (5,0);
    \coordinate (G) at (0,-5);
    \coordinate (C) at (1,4);
    \coordinate (K) at (1,-4);
    \coordinate (D) at (9,4);
    \coordinate (E) at (9,-4);
    \fill[gray!35] (C) -- (Rt) -- (K) -- cycle;
    \draw (A) -- (B) -- (Rt) -- (G) -- cycle;
    \draw (C) -- (D) -- (E) -- (K);
    \draw (C) -- (K);
    \node[left] at (A) {$A$};
    \node[above] at (B) {$B$};
    \node[above right] at (C) {$C$};
    \node[above right] at (D) {$D$};
    \node[below right] at (E) {$E$};
    \node[below] at (G) {$G$};
\end{tikzpicture}
\end{minipage}
\par\vspace{0.2cm}
\answerTable{$164$ см$^2$}{$196$ см$^2$}{$172$ см$^2$}{$188$ см$^2$}{$180$ см$^2$}
\vspace{0.3cm}
\noindent
\begin{minipage}[c]{0.55\textwidth}
\zadnum Для виготовлення картонної моделі трафарету $ABCDE$ Тарас узяв прямокутник $ABCE$ зі сторонами $8$ і $12$~см і вирізав із нього прямокутний трикутник $CDE$ із катетами $6$ і $8$~см (див. рисунок). Обчисліть площу трафарету $ABCDE$.
\end{minipage}%
\hfill
\begin{minipage}[c]{0.43\textwidth}
\centering
\begin{tikzpicture}[scale=0.30, thick]
    \coordinate (A) at (0,0);
    \coordinate (B) at (0,8);
    \coordinate (C) at (12,8);
    \coordinate (E) at (12,0);
    \coordinate (D) at (12,2);
    \coordinate (Dp) at (6,8);
    \fill[gray!35] (Dp) -- (C) -- (D) -- cycle;
    \draw (A) -- (B) -- (Dp) -- (D) -- (E) -- cycle;
    \draw[dashed] (Dp) -- (C) -- (D);
    \node[below left] at (A) {$A$};
    \node[above left] at (B) {$B$};
    \node[above] at (Dp) {$D$};
    \node[right] at (D) {$C$};
    \node[below right] at (E) {$E$};
\end{tikzpicture}
\end{minipage}
\par\vspace{0.2cm}
\answerTable{$72$ см$^2$}{$84$ см$^2$}{$78$ см$^2$}{$90$ см$^2$}{$96$ см$^2$}
\vspace{0.3cm}
\newpage
\taskBlock{Планіметрія: площа складеної фігури --- сесія 30.05, завдання №15}
\noindent
\begin{minipage}[c]{0.51\textwidth}
\zadnum На рисунку зображено схему поверхні декоративного басейну. Воно має форму прямокутника $ABCD$, з кутів якого вирізано чотири однакові чверті кола радіусом $5$~м, а до сторін $BC$ та $AD$ долучено два однакові півкола радіусом $5$~м. Відстань між основами цих півкіл дорівнює $10$~м, а ширина прямокутної частини між дугами вирізів та півкіл становить $10$~м. Знайдіть площу (у м$^2$) поверхні цього басейну.
\end{minipage}%
\hfill
\begin{minipage}[c]{0.48\textwidth}
\centering
\begin{tikzpicture}[scale=0.18, thick]
    \draw[dashed, gray] (0,0) rectangle (40,20);
    \node[below left] at (0,0) {$A$};
    \node[above left] at (0,20) {$B$};
    \node[above right] at (40,20) {$C$};
    \node[below right] at (40,0) {$D$};
    \draw[fill=mainGreen!35, draw=mainGreen!80!black, thick]
        (40,5)
        -- (40,15)
        arc (270:180:5)
        -- (25,20)
        arc (0:180:5)
        -- (5,20)
        arc (360:270:5)
        -- (0,15)
        -- (0,5)
        arc (90:0:5)
        -- (15,0)
        arc (180:360:5)
        -- (35,0)
        arc (180:90:5)
        -- cycle;
    \draw[<->, >=stealth] (-2, 5) -- (-2, 15) node[midway, fill=white, inner sep=1pt] {$10$};
    \draw[<->, >=stealth] (5, 22) -- (15, 22) node[midway, fill=white, inner sep=1pt] {$10$};
    \draw[<->, >=stealth] (25, 22) -- (35, 22) node[midway, fill=white, inner sep=1pt] {$10$};
    \draw[->, >=stealth, gray!90] (20,20) -- (23.5, 23.5) node[right, text=black, font=\scriptsize] {$R=5$};
    \fill[gray] (20,20) circle [radius=0.3];
    \draw[->, >=stealth, gray!90] (0,20) -- (3.5, 16.5) node[right, text=black, font=\scriptsize] {$R=5$};
    \fill[gray] (0,20) circle [radius=0.3];
\end{tikzpicture}
\end{minipage}
\vspace{0.2cm}
\answerTable{$500$}{$50$}{$1000$}{$800$}{$100$}
\hypertarget{task:37}{}\par\vspace{0.05cm}\noindent\hyperlink{sol:37}{\fbox{\footnotesize розв'язок $\rightarrow$}}\par\vspace{0.2cm}
\noindent
\begin{minipage}[c]{0.51\textwidth}
\zadnum На рисунку зображено схему поверхні декоративного басейну. Воно має форму прямокутника $ABCD$, з кутів якого вирізано чотири однакові чверті кола радіусом $4$~м, а до сторін $BC$ та $AD$ долучено два однакові півкола радіусом $4$~м. Відстань між основами цих півкіл дорівнює $8$~м, а ширина прямокутної частини між дугами вирізів та півкіл становить $8$~м. Знайдіть площу (у м$^2$) поверхні цього басейну.
\end{minipage}%
\hfill
\begin{minipage}[c]{0.48\textwidth}
\centering
\begin{tikzpicture}[scale=0.225, thick]
    \draw[dashed, gray] (0,0) rectangle (32,16);
    \node[below left] at (0,0) {$A$};
    \node[above left] at (0,16) {$B$};
    \node[above right] at (32,16) {$C$};
    \node[below right] at (32,0) {$D$};
    \draw[fill=mainGreen!35, draw=mainGreen!80!black, thick]
        (32,4)
        -- (32,12)
        arc (270:180:4)
        -- (20,16)
        arc (0:180:4)
        -- (4,16)
        arc (360:270:4)
        -- (0,12)
        -- (0,4)
        arc (90:0:4)
        -- (12,0)
        arc (180:360:4)
        -- (28,0)
        arc (180:90:4)
        -- cycle;
    \draw[<->, >=stealth] (-1.6, 4) -- (-1.6, 12) node[midway, fill=white, inner sep=1pt] {$8$};
    \draw[<->, >=stealth] (4, 17.6) -- (12, 17.6) node[midway, fill=white, inner sep=1pt] {$8$};
    \draw[<->, >=stealth] (20, 17.6) -- (28, 17.6) node[midway, fill=white, inner sep=1pt] {$8$};
    \draw[->, >=stealth, gray!90] (16,16) -- (18.8, 18.8) node[right, text=black, font=\scriptsize] {$R=4$};
    \fill[gray] (16,16) circle [radius=0.25];
    \draw[->, >=stealth, gray!90] (0,16) -- (2.8, 13.2) node[right, text=black, font=\scriptsize] {$R=4$};
    \fill[gray] (0,16) circle [radius=0.25];
\end{tikzpicture}
\end{minipage}
\vspace{0.2cm}
\answerTable{$320$}{$256$}{$512$}{$640$}{$1024$}
\noindent
\begin{minipage}[c]{0.51\textwidth}
\zadnum На рисунку зображено схему садової ділянки $ABCD$, яка має форму прямокутника зі сторонами $30$~м і $20$~м. У двох протилежних кутах ($A$ і $C$) ділянки вирізано квадратні клумби зі стороною $4$~м, а у двох інших кутах ($B$ і $D$) до ділянки добудовано квадратні тераси зі стороною $4$~м. Знайдіть площу (у м$^2$) цієї ділянки.
\end{minipage}%
\hfill
\begin{minipage}[c]{0.48\textwidth}
\centering
\begin{tikzpicture}[scale=0.13, thick]
    % main figure outline (rectangle minus cutout corners at A,C plus added squares at B,D)
    \draw[fill=mainGreen!35, draw=mainGreen!80!black, thick]
        (4,0) -- (30,0)
        -- (30,-4) -- (34,-4) -- (34,0) -- (30,0)
        -- (30,16) -- (26,16) -- (26,20) -- (0,20)
        -- (-4,20) -- (-4,16) -- (0,16) -- (0,4)
        -- (4,4) -- cycle;
    \draw[dashed, gray!70] (0,0) rectangle (30,20);
    \node[below, font=\small] at (2,0) {$A$};
    \node[above, font=\small] at (2,20) {$B$};
    \node[above, font=\small] at (28,20) {$C$};
    \node[below, font=\small] at (28,0) {$D$};
    \node[font=\tiny] at (15, -2.5) {$30$};
    \node[font=\tiny, rotate=90] at (-2, 10) {$20$};
\end{tikzpicture}
\end{minipage}
\vspace{0.2cm}
\answerTable{$480$}{$568$}{$720$}{$600$}{$632$}
\newpage
\taskBlock{Площа кільця (формула) --- сесія 5.06, завдання №3}
\noindent
\begin{minipage}[c]{0.6\textwidth}
\zadnum Плоска сітчаста секція сушарки має вигляд частини круга, яку обмежено колами радіусів $R$ см і $3$ см, $R>3$ (див. рисунок). За якою формулою можна обчислити площу (у см$^2$) $S$ такої секції?
\end{minipage}%
\hfill
\begin{minipage}[c]{0.38\textwidth}
\centering
\begin{tikzpicture}[scale=0.5, thick]
    \filldraw[fill=gray!15, draw=cyan!60!black, thick] (0,0) circle (2);
    \filldraw[fill=white, draw=cyan!60!black, thick] (0,0) circle (0.6);
\end{tikzpicture}
\end{minipage}
\par\vspace{0.1cm}
\begin{flushleft}
\textbf{А}\ \ $S = 2\pi R - 6\pi$ \qquad \textbf{Б}\ \ $S = \pi R^2 - 6\pi$ \qquad \textbf{В}\ \ $S = \pi(R-3)^2$\\[0.15cm]
\textbf{Г}\ \ $S = \pi R^2 + 9\pi$ \qquad \textbf{Д}\ \ $S = \pi R^2 - 9\pi$
\end{flushleft}
\hypertarget{task:40}{}\par\vspace{0.05cm}\noindent\hyperlink{sol:40}{\fbox{\footnotesize розв'язок $\rightarrow$}}\par\vspace{0.2cm}
\vspace{0.2cm}
\noindent
\begin{minipage}[c]{0.6\textwidth}
\zadnum Металева прокладка має вигляд частини круга, яку обмежено колами радіусів $R$ см і $4$ см, $R>4$ (див. рисунок). За якою формулою можна обчислити площу (у см$^2$) $S$ такої прокладки?
\end{minipage}%
\hfill
\begin{minipage}[c]{0.38\textwidth}
\centering
\begin{tikzpicture}[scale=0.5, thick]
    \filldraw[fill=gray!15, draw=cyan!60!black, thick] (0,0) circle (2);
    \filldraw[fill=white, draw=cyan!60!black, thick] (0,0) circle (0.8);
\end{tikzpicture}
\end{minipage}
\par\vspace{0.1cm}
\begin{flushleft}
\textbf{А}\ \ $S = 2\pi R - 8\pi$ \qquad \textbf{Б}\ \ $S = \pi R^2 - 8\pi$ \qquad \textbf{В}\ \ $S = \pi(R-4)^2$\\[0.15cm]
\textbf{Г}\ \ $S = \pi R^2 - 16\pi$ \qquad \textbf{Д}\ \ $S = \pi R^2 + 16\pi$
\end{flushleft}
\vspace{0.2cm}
\noindent
\begin{minipage}[c]{0.6\textwidth}
\zadnum Декоративна металева пластина має вигляд півкільця, обмеженого двома півколами радіусів $R$ см та $3$ см, $R > 3$ (див. рисунок). За якою формулою можна обчислити площу (у см$^2$) $S$ такої пластини?
\end{minipage}%
\hfill
\begin{minipage}[c]{0.38\textwidth}
\centering
\begin{tikzpicture}[scale=0.5, thick]
    \filldraw[fill=gray!15, draw=cyan!60!black, thick] (-2,0) arc (180:0:2) -- (0.6,0) arc (0:180:0.6) -- cycle;
    \draw[cyan!60!black, thick] (-2,0) -- (2,0);
\end{tikzpicture}
\end{minipage}
\par\vspace{0.1cm}
\begin{flushleft}
\textbf{А}\ \ $S = \dfrac{\pi R^2}{2} - \dfrac{9\pi}{2}$ \qquad \textbf{Б}\ \ $S = \pi R^2 - 9\pi$ \qquad \textbf{В}\ \ $S = \dfrac{\pi(R-3)^2}{2}$\\[0.15cm]
\textbf{Г}\ \ $S = \dfrac{\pi R^2}{2} - 9\pi$ \qquad \textbf{Д}\ \ $S = \dfrac{\pi R^2}{2} + \dfrac{9\pi}{2}$
\end{flushleft}
\vspace{0.2cm}
\newpage
\subTypeTitle{Твердження та властивості геометричних фігур}
\taskBlock{Істинність тверджень про коло --- сесія 23.05, завдання №10}
\zadtask{Які з наведених тверджень є істинними?
\vspace{0.15cm}
\noindent\hspace{1em}I. Пряма, що проходить через центр кола і лежить з ним в одній площині, має з ним дві спільні точки. \\
\noindent\hspace{1em}II. Діаметр кола, перпендикулярний до хорди, проходить через середину цієї хорди. \\
\noindent\hspace{1em}III. У колі можна провести два діаметри, які не перетинаються.}
\answerTable{лише I і II}{лише II}{лише I}{лише II і III}{усі три}
\hypertarget{task:43}{}\par\vspace{0.05cm}\noindent\hyperlink{sol:43}{\fbox{\footnotesize розв'язок $\rightarrow$}}\par\vspace{0.2cm}
\vspace{0.3cm}
\zadtask{Які з наведених тверджень є істинними?
\vspace{0.15cm}
\noindent\hspace{1em}I. Дотична до кола перпендикулярна до радіуса, проведеного в точку дотику. \\
\noindent\hspace{1em}II. Через будь-які три точки площини можна провести коло. \\
\noindent\hspace{1em}III. Вписаний кут, що спирається на діаметр кола, є прямим.}
\answerTable{лише I}{лише II і III}{лише I і II}{лише I і III}{усі три}
\vspace{0.3cm}
\zadtask{Які з наведених тверджень є істинними?
\vspace{0.15cm}
\noindent\hspace{1em}I. Центральний кут удвічі більший за вписаний кут, що спирається на ту саму дугу. \\
\noindent\hspace{1em}II. Хорди кола, рівновіддалені від його центра, рівні між собою. \\
\noindent\hspace{1em}III. Відрізки дотичних, проведених до кола з однієї точки поза колом, рівні між собою.}
\answerTable{лише I}{лише II}{лише I і III}{лише II і III}{усі три}
\vspace{0.3cm}
\newpage
\taskBlock{Планіметрія: твердження про трикутники --- сесія 30.05, завдання №9}
\zadtask{Розгляньте наведені нижче твердження.
\vspace{0.2cm}
\noindent
\makebox[1.8em][l]{\textbf{I}}\, Найбільший кут прямокутного трикутника дорівнює половині розгорнутого кута.\par\vspace{0.1cm}
\makebox[1.8em][l]{\textbf{II}}\, Сума двох кутів рівностороннього трикутника дорівнює $120^\circ$.\par\vspace{0.1cm}
\makebox[1.8em][l]{\textbf{III}}\, Усі гострі кути гострокутного трикутника більші за $60^\circ$.\par\vspace{0.25cm}
\noindent Які з тверджень є \textit{правильними}?
\par\vspace{0.2cm}
\begin{flushleft}
\textbf{А} \quad лише I \\[0.1cm]
\textbf{Б} \quad лише II \\[0.1cm]
\textbf{В} \quad лише I та II \\[0.1cm]
\textbf{Г} \quad лише II та III \\[0.1cm]
\textbf{Д} \quad I, II та III
\end{flushleft}}
\hypertarget{task:46}{}\par\vspace{0.05cm}\noindent\hyperlink{sol:46}{\fbox{\footnotesize розв'язок $\rightarrow$}}\par\vspace{0.2cm}
\vspace{0.3cm}
\zadtask{Розгляньте наведені нижче твердження.
\vspace{0.2cm}
\noindent
\makebox[1.8em][l]{\textbf{I}}\, Сума гострих кутів прямокутного трикутника дорівнює $90^\circ$.\par\vspace{0.1cm}
\makebox[1.8em][l]{\textbf{II}}\, У тупокутному трикутнику два кути можуть бути тупими.\par\vspace{0.1cm}
\makebox[1.8em][l]{\textbf{III}}\, Кожен кут рівностороннього трикутника дорівнює третині розгорнутого кута.\par\vspace{0.25cm}
\noindent Які з тверджень є \textit{правильними}?
\par\vspace{0.2cm}
\begin{flushleft}
\textbf{А} \quad лише I \\[0.1cm]
\textbf{Б} \quad лише III \\[0.1cm]
\textbf{В} \quad лише I та III \\[0.1cm]
\textbf{Г} \quad лише II та III \\[0.1cm]
\textbf{Д} \quad I, II та III
\end{flushleft}}
\vspace{0.3cm}
\zadtask{Розгляньте наведені нижче твердження.
\vspace{0.2cm}
\noindent
\makebox[1.8em][l]{\textbf{I}}\, Зовнішній кут трикутника дорівнює сумі двох внутрішніх кутів, не суміжних з ним.\par\vspace{0.1cm}
\makebox[1.8em][l]{\textbf{II}}\, У рівнобедреному трикутнику кути при основі рівні.\par\vspace{0.1cm}
\makebox[1.8em][l]{\textbf{III}}\, У будь-якому трикутнику найменший кут менший за $90^\circ$.\par\vspace{0.25cm}
\noindent Які з тверджень є \textit{правильними}?
\par\vspace{0.2cm}
\begin{flushleft}
\textbf{А} \quad лише I \\[0.1cm]
\textbf{Б} \quad лише II \\[0.1cm]
\textbf{В} \quad лише I та II \\[0.1cm]
\textbf{Г} \quad лише II та III \\[0.1cm]
\textbf{Д} \quad I, II та III
\end{flushleft}}
\vspace{0.3cm}
\newpage
\taskBlock{Рівнобедрений трикутник, властивості бісектриси (висоти, медіани) --- сесія 1.06, завдання №9}
\zadtask{У рівнобедреному трикутнику $ABC$ ($AB = BC$) $BK$~--- бісектриса кута $B$, точка $K$ належить стороні $AC$. Які з наведених тверджень є \textit{правильними}?
\vspace{0.2cm}
\noindent
\makebox[1.8em][l]{\textbf{I}}\, $AB + BC = AC$.\par\vspace{0.1cm}
\makebox[1.8em][l]{\textbf{II}}\, $BK \perp AC$.\par\vspace{0.1cm}
\makebox[1.8em][l]{\textbf{III}}\, $AK = KC$.}
\answerTable{лише I}{лише I та II}{лише I та III}{лише II та III}{I, II та III}
\hypertarget{task:49}{}\par\vspace{0.05cm}\noindent\hyperlink{sol:49}{\fbox{\footnotesize розв'язок $\rightarrow$}}\par\vspace{0.2cm}
\zadtask{У рівнобедреному трикутнику $ABC$ ($AB = BC$) $BM$~--- медіана, точка $M$ належить стороні $AC$. Які з наведених тверджень є \textit{правильними}?
\vspace{0.2cm}
\noindent
\makebox[1.8em][l]{\textbf{I}}\, $BM$~--- бісектриса кута $B$.\par\vspace{0.1cm}
\makebox[1.8em][l]{\textbf{II}}\, $AM = AB$.\par\vspace{0.1cm}
\makebox[1.8em][l]{\textbf{III}}\, $BM \perp AC$.}
\answerTable{лише I}{лише II}{лише I та III}{лише II та III}{I, II та III}
\zadtask{У рівнобедреному трикутнику $ABC$ ($AB = BC$) $BH$~--- висота, точка $H$ належить стороні $AC$. Які з наведених тверджень є \textit{правильними}?
\vspace{0.2cm}
\noindent
\makebox[1.8em][l]{\textbf{I}}\, $AH = HC$.\par\vspace{0.1cm}
\makebox[1.8em][l]{\textbf{II}}\, $\angle ABH = \angle CBH$.\par\vspace{0.1cm}
\makebox[1.8em][l]{\textbf{III}}\, $AB = AC$.}
\answerTable{лише I}{лише I та II}{лише I та III}{лише II та III}{I, II та III}
\newpage
\taskBlock{Твердження про елементи трикутника та чотирикутника --- сесія 2.06, завдання №9}
\zadtask{Які з наведених тверджень є \textit{правильними}?
\vspace{0.1cm}
\begin{enumerate}[label=\textbf{\Roman*}, align=left, leftmargin=2.8em, labelsep=0.6em, itemsep=0.12cm, topsep=0.1cm]
\item У правильному трикутнику всі медіани рівні.
\item Медіана трикутника ділить його на два трикутники з рівними периметрами.
\item Медіана трикутника ділить його на два трикутники з рівними площами.
\end{enumerate}
\vspace{0.15cm}}
\answerTable{лише I}{лише II}{лише I та III}{лише II та III}{I, II та III}
\hypertarget{task:52}{}\par\vspace{0.05cm}\noindent\hyperlink{sol:52}{\fbox{\footnotesize розв'язок $\rightarrow$}}\par\vspace{0.2cm}
\zadtask{Які з наведених тверджень є \textit{правильними}?
\vspace{0.1cm}
\begin{enumerate}[label=\textbf{\Roman*}, align=left, leftmargin=2.8em, labelsep=0.6em, itemsep=0.12cm, topsep=0.1cm]
\item У рівносторонньому трикутнику всі висоти рівні.
\item Висота трикутника завжди лежить усередині трикутника.
\item Бісектриса будь-якого трикутника ділить протилежну сторону навпіл.
\end{enumerate}
\vspace{0.15cm}}
\answerTable{лише I}{лише II}{лише I та II}{лише I та III}{I, II та III}
\zadtask{Які з наведених тверджень є \textit{правильними}?
\vspace{0.1cm}
\begin{enumerate}[label=\textbf{\Roman*}, align=left, leftmargin=2.8em, labelsep=0.6em, itemsep=0.12cm, topsep=0.1cm]
\item Діагоналі ромба перпендикулярні.
\item Діагоналі ромба рівні.
\item Діагоналі ромба є бісектрисами його кутів.
\end{enumerate}
\vspace{0.15cm}}
\answerTable{лише I}{лише II}{лише I та III}{лише II та III}{I, II та III}
\newpage
\taskBlock{Планіметрія: властивості кола --- сесія 3.06, завдання №9}
\zadnum Які з наведених тверджень є \textit{правильними}?
\vspace{0.1cm}
\begin{enumerate}[label=\textbf{\Roman*.}, align=left, leftmargin=2.6em, labelsep=0.5em, itemsep=0.1cm, topsep=0.08cm]
\item Будь-яка хорда кола більша за його радіус.
\item Кінці діаметра кола ділять коло навпіл.
\item Рівні хорди кола стягують рівні дуги.
\end{enumerate}
\vspace{0.12cm}
\answerTable{лише II}{лише I та II}{лише III}{лише II та III}{I, II та III}
\hypertarget{task:55}{}\par\vspace{0.05cm}\noindent\hyperlink{sol:55}{\fbox{\footnotesize розв'язок $\rightarrow$}}\par\vspace{0.2cm}
\zadnum Які з наведених тверджень є \textit{правильними}?
\vspace{0.1cm}
\begin{enumerate}[label=\textbf{\Roman*.}, align=left, leftmargin=2.6em, labelsep=0.5em, itemsep=0.1cm, topsep=0.08cm]
\item Діаметр є найбільшою хордою кола.
\item Будь-які два радіуси одного кола рівні між собою.
\item Будь-який діаметр кола ділить будь-яку його хорду навпіл.
\end{enumerate}
\vspace{0.12cm}
\answerTable{лише I}{лише I та II}{лише II та III}{лише III}{I, II та III}
\zadnum Які з наведених тверджень є \textit{неправильними}?
\vspace{0.1cm}
\begin{enumerate}[label=\textbf{\Roman*.}, align=left, leftmargin=2.6em, labelsep=0.5em, itemsep=0.1cm, topsep=0.08cm]
\item Через будь-які три точки площини можна провести коло.
\item Центр кола, описаного навколо трикутника, завжди лежить усередині трикутника.
\item Серединний перпендикуляр до хорди проходить через центр кола.
\end{enumerate}
\vspace{0.12cm}
\answerTable{лише I}{лише I та II}{лише II}{лише II та III}{I, II та III}
\newpage
\taskBlock{Геометричні твердження (коло і трикутник) --- сесія 4.06, завдання №9}
\zadtask{Які з наведених тверджень є правильними?
\vspace{0.1cm}
\par\noindent\hangindent=2.2em\textbf{I.}\ \ Коло, уписане в рівнобедрений трикутник, дотикається до середини його основи.
\par\noindent\hangindent=2.2em\textbf{II.}\ \ Центр кола, описаного навколо тупокутного трикутника, лежить усередині трикутника.
\par\noindent\hangindent=2.2em\textbf{III.}\ \ Коло, побудоване на гіпотенузі прямокутного трикутника як на діаметрі, проходить через вершину прямого кута.
\vspace{0.15cm}
\answerTable{лише I}{лише I та II}{I, II та III}{лише II та III}{лише I та III}}
\hypertarget{task:58}{}\par\vspace{0.05cm}\noindent\hyperlink{sol:58}{\fbox{\footnotesize розв'язок $\rightarrow$}}\par\vspace{0.2cm}
\zadtask{Які з наведених тверджень є правильними?
\vspace{0.1cm}
\par\noindent\hangindent=2.2em\textbf{I.}\ \ Центр кола, описаного навколо прямокутного трикутника, лежить на середині гіпотенузи.
\par\noindent\hangindent=2.2em\textbf{II.}\ \ Коло, уписане в трикутник, дотикається до всіх трьох його сторін.
\par\noindent\hangindent=2.2em\textbf{III.}\ \ Центр кола, описаного навколо гострокутного трикутника, лежить поза трикутником.
\vspace{0.15cm}
\answerTable{лише I}{I, II та III}{лише I та III}{лише II та III}{лише I та II}}
\zadtask{Які з наведених тверджень є правильними?
\vspace{0.1cm}
\par\noindent\hangindent=2.2em\textbf{I.}\ \ Центр кола, уписаного в трикутник, є точкою перетину його бісектрис.
\par\noindent\hangindent=2.2em\textbf{II.}\ \ Радіус кола, описаного навколо рівностороннього трикутника, удвічі більший за радіус кола, уписаного в цей трикутник.
\par\noindent\hangindent=2.2em\textbf{III.}\ \ Центр кола, описаного навколо будь-якого трикутника, завжди лежить усередині трикутника.
\vspace{0.15cm}
\answerTable{лише I}{лише I та II}{лише I та III}{лише II та III}{I, II та III}}
\newpage
\taskBlock{Середня лінія трикутника (твердження) --- сесія 5.06, завдання №9}
\zadnum На сторонах $AB$ та $BC$ довільного трикутника $ABC$ вибрано точки $M$ та $N$ так, що $AM = MB$, $CN = NB$. Які з наведених тверджень є \textit{правильними}?
\vspace{0.1cm}
\begin{enumerate}[label=\textbf{\Roman*.}, align=left, leftmargin=2.6em, labelsep=0.5em, itemsep=0.08cm, topsep=0.06cm]
\item Чотирикутник $AMNC$ є трапецією.
\item Трикутник $MBN$ рівнобедрений.
\item Периметр трикутника $MBN$ дорівнює половині периметра трикутника $ABC$.
\end{enumerate}
\vspace{0.1cm}
\answerTable{лише II}{лише I та II}{лише I та III}{лише II та III}{I, II та III}
\hypertarget{task:61}{}\par\vspace{0.05cm}\noindent\hyperlink{sol:61}{\fbox{\footnotesize розв'язок $\rightarrow$}}\par\vspace{0.2cm}
\zadnum На сторонах $DE$ та $EF$ довільного трикутника $DEF$ вибрано точки $P$ та $Q$ так, що $DP = PE$, $FQ = QE$. Які з наведених тверджень є \textit{правильними}?
\vspace{0.1cm}
\begin{enumerate}[label=\textbf{\Roman*.}, align=left, leftmargin=2.6em, labelsep=0.5em, itemsep=0.08cm, topsep=0.06cm]
\item Відрізок $PQ$ паралельний стороні $DF$.
\item Довжина відрізка $PQ$ удвічі менша за довжину сторони $DF$.
\item Трикутник $PEQ$ подібний до трикутника $DEF$ з коефіцієнтом подібності $\dfrac{1}{3}$.
\end{enumerate}
\vspace{0.1cm}
\answerTable{лише I}{лише I та II}{лише I та III}{лише II та III}{I, II та III}
\zadnum Точки $M$ і $N$~--- середини сторін $AB$ та $BC$ довільного трикутника $ABC$. Які з наведених тверджень є \textit{правильними}?
\vspace{0.1cm}
\begin{enumerate}[label=\textbf{\Roman*.}, align=left, leftmargin=2.6em, labelsep=0.5em, itemsep=0.08cm, topsep=0.06cm]
\item Площа трикутника $MBN$ становить половину площі трикутника $ABC$.
\item Довжина відрізка $MN$ удвічі менша за довжину сторони $AC$.
\item Відрізок $MN$ паралельний стороні $AC$.
\end{enumerate}
\vspace{0.1cm}
\answerTable{лише I}{лише I та II}{лише I та III}{лише II та III}{I, II та III}
\newpage
\typeTitle{Завдання на встановлення відповідності}
\subTypeTitle{Властивості геометричних фігур (відповідність)}
\taskBlock{Трапеція з висотою-бісектрисою (відповідність) --- сесія 23.05, завдання №16}
\zadtask{На рисунку зображено трапецію $ABCD$, у якій $\angle CDA = 30^\circ$. Висота $CK$, довжина якої дорівнює $2$, є бісектрисою кута $ACD$. \nmtyear{2026}
\vspace{0.3cm}
\begin{center}
\begin{tikzpicture}[scale=0.9, thick]
    \coordinate (A) at (0,0);
    \coordinate (K) at (4,0);
    \coordinate (D) at (7.46,0);
    \coordinate (C) at (4,2);
    \coordinate (B) at (0,2);
    \draw (A) -- (B) -- (C) -- (D) -- cycle;
    \draw (C) -- (K);
    \draw[dashed] (A) -- (C);
    \draw (0,0) -- (0,0.3) -- (0.3,0.3) -- (0.3,0);
    \draw (4,0) -- (4,0.3) -- (4.3,0.3) -- (4.3,0);
    \pic [draw, angle radius=0.6cm, "$30^\circ$"{above left, shift={(-0.28,-0.15)}}] {angle = C--D--K};
    \pic [draw, angle radius=0.45cm] {angle = K--C--D};
    \pic [draw, angle radius=0.55cm] {angle = A--C--K};
    \node[below left] at (A) {$A$};
    \node[above left] at (B) {$B$};
    \node[above=2pt] at (C) {$C$};
    \node[below right] at (D) {$D$};
    \node[below=2pt] at (K) {$K$};
\end{tikzpicture}
\end{center}
\vspace{0.3cm}
\noindent Установіть відповідність між величиною (1--3) та її значенням (А--Д).
\vspace{0.3cm}
\noindent
\begin{minipage}[t]{0.52\textwidth}
\textit{Величина} \par \vspace{0.2cm}
\textbf{1} \ \ довжина відрізка $CD$\\[0.25cm]
\textbf{2} \ \ довжина відрізка $BC$\\[0.25cm]
\textbf{3} \ \ відстань від точки $B$ до прямої $AC$
\end{minipage}
\begin{minipage}[t]{0.28\textwidth}
\textit{Значення} \par \vspace{0.2cm}
\textbf{А} \ \ $2\sqrt{3}$\\[0.15cm]
\textbf{Б} \ \ $4$\\[0.15cm]
\textbf{В} \ \ $1$\\[0.15cm]
\textbf{Г} \ \ $\sqrt{3}$\\[0.15cm]
\textbf{Д} \ \ $6\sqrt{3}$
\end{minipage}
\hfill
\begin{minipage}[t]{0.22\textwidth}
\vspace{-0.2cm}
\begin{flushright}\matchingGrid\end{flushright}
\end{minipage}
}
\hypertarget{task:64}{}\par\vspace{0.05cm}\noindent\hyperlink{sol:64}{\fbox{\footnotesize розв'язок $\rightarrow$}}\par\vspace{0.2cm}
\vspace{0.5cm}
\zadtask{На рисунку зображено трапецію $ABCD$, у якій $\angle CDA = 45^\circ$. Висота $CK$, довжина якої дорівнює $3$, є бісектрисою кута $ACD$.
\vspace{0.3cm}
\begin{center}
\begin{tikzpicture}[scale=0.9, thick]
    \coordinate (A) at (0,0);
    \coordinate (K) at (3,0);
    \coordinate (D) at (6,0);
    \coordinate (C) at (3,3);
    \coordinate (B) at (0,3);
    \draw (A) -- (B) -- (C) -- (D) -- cycle;
    \draw (C) -- (K);
    \draw[dashed] (A) -- (C);
    \draw (0,0) -- (0,0.3) -- (0.3,0.3) -- (0.3,0);
    \draw (3,0) -- (3,0.3) -- (3.3,0.3) -- (3.3,0);
    \pic [draw, angle radius=0.6cm, "$45^\circ$"{above left, shift={(-0.28,-0.15)}}] {angle = C--D--K};
    \pic [draw, angle radius=0.45cm] {angle = K--C--D};
    \pic [draw, angle radius=0.55cm] {angle = A--C--K};
    \node[below left] at (A) {$A$};
    \node[above left] at (B) {$B$};
    \node[above=2pt] at (C) {$C$};
    \node[below right] at (D) {$D$};
    \node[below=2pt] at (K) {$K$};
\end{tikzpicture}
\end{center}
\vspace{0.3cm}
\noindent Установіть відповідність між величиною (1--3) та її значенням (А--Д).
\vspace{0.3cm}
\noindent
\begin{minipage}[t]{0.52\textwidth}
\textit{Величина} \par \vspace{0.2cm}
\textbf{1} \ \ довжина відрізка $CD$\\[0.25cm]
\textbf{2} \ \ довжина відрізка $BC$\\[0.25cm]
\textbf{3} \ \ відстань від точки $B$ до прямої $AC$
\end{minipage}
\begin{minipage}[t]{0.28\textwidth}
\textit{Значення} \par \vspace{0.2cm}
\textbf{А} \ \ $3$\\[0.15cm]
\textbf{Б} \ \ $\dfrac{3\sqrt{2}}{2}$\\[0.15cm]
\textbf{В} \ \ $6$\\[0.15cm]
\textbf{Г} \ \ $3\sqrt{2}$\\[0.15cm]
\textbf{Д} \ \ $2\sqrt{3}$
\end{minipage}
\hfill
\begin{minipage}[t]{0.22\textwidth}
\vspace{-0.2cm}
\begin{flushright}\matchingGrid\end{flushright}
\end{minipage}
}
\vspace{0.5cm}
\zadtask{У прямокутнику $ABCD$ на стороні $AD$ вибрано точку $K$ так, що $AK : KD = 2 : 3$ (див. рисунок). $BK = 10$~см, $CD = 8$~см. Узгодьте початок речення (1--3) із його закінченням (А--Д) так, щоб утворилося правильне твердження.
\vspace{0.3cm}
\begin{center}
\begin{tikzpicture}[scale=0.75, thick]
    \coordinate (A) at (0,0);
    \coordinate (D) at (7.5,0);
    \coordinate (C) at (7.5,4);
    \coordinate (B) at (0,4);
    \coordinate (K) at (3,0);
    \draw (A) -- (B) -- (C) -- (D) -- cycle;
    \draw (B) -- (K);
    \node[below left] at (A) {$A$};
    \node[above left] at (B) {$B$};
    \node[above right] at (C) {$C$};
    \node[below right] at (D) {$D$};
    \node[below] at (K) {$K$};
\end{tikzpicture}
\end{center}
\vspace{0.2cm}
\noindent
\begin{minipage}[t]{0.50\textwidth}
\textit{Початок речення} \par \vspace{0.15cm}
\textbf{1} \ \ Довжина відрізка $AK$ дорівнює\\[0.2cm]
\textbf{2} \ \ Середня лінія трапеції $BKDC$ дорівнює\\[0.2cm]
\textbf{3} \ \ Діаметр кола, описаного навколо прямокутника $ABCD$, дорівнює
\end{minipage}
\begin{minipage}[t]{0.28\textwidth}
\textit{Закінчення речення} \par \vspace{0.15cm}
\textbf{А} \ \ $6$ см\\[0.1cm]
\textbf{Б} \ \ $9$ см\\[0.1cm]
\textbf{В} \ \ $12$ см\\[0.1cm]
\textbf{Г} \ \ $15$ см\\[0.1cm]
\textbf{Д} \ \ $17$ см
\end{minipage}
\hfill
\begin{minipage}[t]{0.20\textwidth}
\vspace{-0.2cm}
\begin{flushright}\matchingGrid\end{flushright}
\end{minipage}
}
\newpage
\taskBlock{Планіметрія: квадрат і трикутник (відповідність) --- сесія 30.05, завдання №16}
\zadtask{На рисунку зображено квадрат $ABCD$ та рівнобедрений прямокутний трикутник $KBC$ (прямий кут при вершині $K$), які мають спільну сторону $BC$. Точка $M$~--- середина сторони $CD$. Довжина ламаної $BADM$ дорівнює $20$. Установіть відповідність між величинами (1--3) та їхніми значеннями (А--Д).
\vspace{0.3cm}
\begin{center}
\begin{tikzpicture}[scale=0.75, thick]
    \coordinate (A) at (0,0);
    \coordinate (B) at (0,4);
    \coordinate (C) at (4,4);
    \coordinate (D) at (4,0);
    \coordinate (K) at (2,6);
    \coordinate (Mid) at (4,2);
    \draw[thick] (A) -- (B) -- (C) -- (D) -- cycle;
    \draw[thick] (B) -- (K) -- (C);
    \fill (Mid) circle (2pt);
    \node[below left] at (A) {$A$};
    \node[above left] at (B) {$B$};
    \node[above right] at (C) {$C$};
    \node[below right] at (D) {$D$};
    \node[above] at (K) {$K$};
    \node[right] at (Mid) {$M$};
\end{tikzpicture}
\end{center}
\vspace{0.2cm}
\noindent
\begin{minipage}[t]{0.45\textwidth}
\textit{Величина}\par\vspace{0.15cm}
\textbf{1} \quad довжина сторони квадрата\\[0.25cm]
\textbf{2} \quad відстань від точки $K$ до прямої $AD$\\[0.25cm]
\textbf{3} \quad довжина відрізка $KM$
\end{minipage}
\hfill
\begin{minipage}[t]{0.3\textwidth}
\textit{Значення}\par\vspace{0.15cm}
\textbf{А} \quad $12$\\[0.15cm]
\textbf{Б} \quad $10$\\[0.15cm]
\textbf{В} \quad $4\sqrt{3}$\\[0.15cm]
\textbf{Г} \quad $4\sqrt{5}$\\[0.15cm]
\textbf{Д} \quad $8$
\end{minipage}
\hfill
\begin{minipage}[t]{0.2\textwidth}
\vspace{0.4cm}
\begin{flushright}\matchingGrid\end{flushright}
\end{minipage}}
\hypertarget{task:67}{}\par\vspace{0.05cm}\noindent\hyperlink{sol:67}{\fbox{\footnotesize розв'язок $\rightarrow$}}\par\vspace{0.2cm}
\vspace{0.7cm}
\zadtask{На рисунку зображено квадрат $ABCD$ та рівнобедрений прямокутний трикутник $KBC$ (прямий кут при вершині $K$), які мають спільну сторону $BC$. Точка $M$~--- середина сторони $CD$. Довжина ламаної $BADM$ дорівнює $30$. Установіть відповідність між величинами (1--3) та їхніми значеннями (А--Д).
\vspace{0.3cm}
\begin{center}
\begin{tikzpicture}[scale=0.75, thick]
    \coordinate (A) at (0,0);
    \coordinate (B) at (0,4);
    \coordinate (C) at (4,4);
    \coordinate (D) at (4,0);
    \coordinate (K) at (2,6);
    \coordinate (Mid) at (4,2);
    \draw[thick] (A) -- (B) -- (C) -- (D) -- cycle;
    \draw[thick] (B) -- (K) -- (C);
    \fill (Mid) circle (2pt);
    \node[below left] at (A) {$A$};
    \node[above left] at (B) {$B$};
    \node[above right] at (C) {$C$};
    \node[below right] at (D) {$D$};
    \node[above] at (K) {$K$};
    \node[right] at (Mid) {$M$};
\end{tikzpicture}
\end{center}
\vspace{0.2cm}
\noindent
\begin{minipage}[t]{0.45\textwidth}
\textit{Величина}\par\vspace{0.15cm}
\textbf{1} \quad довжина сторони квадрата\\[0.25cm]
\textbf{2} \quad відстань від точки $K$ до прямої $AD$\\[0.25cm]
\textbf{3} \quad довжина відрізка $KM$
\end{minipage}
\hfill
\begin{minipage}[t]{0.3\textwidth}
\textit{Значення}\par\vspace{0.15cm}
\textbf{А} \quad $18$\\[0.15cm]
\textbf{Б} \quad $15$\\[0.15cm]
\textbf{В} \quad $6\sqrt{3}$\\[0.15cm]
\textbf{Г} \quad $6\sqrt{5}$\\[0.15cm]
\textbf{Д} \quad $12$
\end{minipage}
\hfill
\begin{minipage}[t]{0.2\textwidth}
\vspace{0.4cm}
\begin{flushright}\matchingGrid\end{flushright}
\end{minipage}}
\vspace{0.7cm}
\zadtask{У ромбі $ABCD$ з вершини $B$ опущено висоту $BH$ на сторону $AD$ (див. рисунок). $AB = 10$~см, $BH = 8$~см. Установіть відповідність між величиною (1--3) та її значенням (А--Д).
\vspace{0.3cm}
\begin{center}
\begin{tikzpicture}[scale=0.4, thick]
    \coordinate (A) at (0,0);
    \coordinate (D) at (10,0);
    \coordinate (B) at (6,8);
    \coordinate (Cc) at (16,8);
    \coordinate (H) at (6,0);
    \draw[thick] (A) -- (B) -- (Cc) -- (D) -- cycle;
    \draw[thick] (B) -- (H);
    \draw (H) ++(0,0.3) -- ++(-0.3,0) -- ++(0,-0.3);
    \fill (H) circle (2pt);
    \node[below left] at (A) {$A$};
    \node[above left] at (B) {$B$};
    \node[above right] at (Cc) {$C$};
    \node[below right] at (D) {$D$};
    \node[below] at (H) {$H$};
\end{tikzpicture}
\end{center}
\vspace{0.2cm}
\noindent
\begin{minipage}[t]{0.46\textwidth}
\textit{Величина} \par \vspace{0.15cm}
\textbf{1} \ \ довжина відрізка $AH$\\[0.25cm]
\textbf{2} \ \ довжина відрізка $HD$\\[0.25cm]
\textbf{3} \ \ площа ромба $ABCD$
\end{minipage}
\begin{minipage}[t]{0.30\textwidth}
\textit{Значення} \par \vspace{0.15cm}
\textbf{А} \ \ $4$ см\\[0.1cm]
\textbf{Б} \ \ $6$ см\\[0.1cm]
\textbf{В} \ \ $10$ см\\[0.1cm]
\textbf{Г} \ \ $64$ см$^2$\\[0.1cm]
\textbf{Д} \ \ $80$ см$^2$
\end{minipage}
\hfill
\begin{minipage}[t]{0.20\textwidth}
\vspace{-0.2cm}
\begin{flushright}\matchingGrid\end{flushright}
\end{minipage}
}
\vspace{0.7cm}
\newpage
\taskBlock{Паралелограм, висота й подібність трикутників (відповідність) --- сесія 1.06, завдання №18}
\zadtask{У паралелограмі $ABCD$ на висоті $BM$ вибрано точку $K$ так, що $BK : KM = 4 : 3$ (див. рисунок). $CK = 17$ см, $BK = 8$ см. Прямі $CK$ і $AD$ перетинаються в точці $N$. Установіть відповідність між відрізком (1--3) та його довжиною (А--Д).
\vspace{0.3cm}
\begin{center}
\begin{tikzpicture}[scale=0.95, thick]
    \coordinate (A) at (0,0);
    \coordinate (M) at (1.3,0);
    \coordinate (D) at (5,0);
    \coordinate (B) at (1.3,2.6);
    \coordinate (C) at (5,2.6);
    \coordinate (N) at (-1.1,0);
    \coordinate (K) at (1.3,1.1);
    \draw[thick] (A) -- (B) -- (C) -- (D) -- cycle;
    \draw[thick] (B) -- (M);
    \draw[thick] (N) -- (C);
    \draw[thick] (N) -- (A);
    \draw[thick] (M)++(0,0.28) -- ++(0.28,0) -- ++(0,-0.28);
    \fill (K) circle (1.6pt);
    \fill (N) circle (1.6pt);
    \node[below] at (A) {$A$};
    \node[below] at (M) {$M$};
    \node[below right] at (D) {$D$};
    \node[above left] at (B) {$B$};
    \node[above right] at (C) {$C$};
    \node[below] at (N) {$N$};
    \node[below right] at (K) {$K$};
\end{tikzpicture}
\end{center}
\vspace{0.2cm}
\noindent
\begin{minipage}[t]{0.42\textwidth}
\textit{Відрізок}\par\vspace{0.2cm}
\textbf{1} \quad $BC$\\[0.25cm]
\textbf{2} \quad $BM$\\[0.25cm]
\textbf{3} \quad $NM$
\end{minipage}
\hfill
\begin{minipage}[t]{0.3\textwidth}
\textit{Довжина відрізка}\par\vspace{0.2cm}
\textbf{А} \quad $11{,}25$ см\\[0.15cm]
\textbf{Б} \quad $12$ см\\[0.15cm]
\textbf{В} \quad $14$ см\\[0.15cm]
\textbf{Г} \quad $15$ см\\[0.15cm]
\textbf{Д} \quad $20$ см
\end{minipage}
\hfill
\begin{minipage}[t]{0.2\textwidth}
\vspace{0.4cm}
\begin{flushright}\matchingGrid\end{flushright}
\end{minipage}}
\hypertarget{task:70}{}\par\vspace{0.05cm}\noindent\hyperlink{sol:70}{\fbox{\footnotesize розв'язок $\rightarrow$}}\par\vspace{0.2cm}
\vspace{0.5cm}
\zadtask{У паралелограмі $ABCD$ на висоті $BM$ вибрано точку $K$ так, що $BK : KM = 3 : 2$ (див. рисунок). $CK = 15$ см, $BK = 9$ см. Прямі $CK$ і $AD$ перетинаються в точці $N$. Установіть відповідність між відрізком (1--3) та його довжиною (А--Д).
\vspace{0.3cm}
\begin{center}
\begin{tikzpicture}[scale=0.95, thick]
    \coordinate (A) at (0,0);
    \coordinate (M) at (1.3,0);
    \coordinate (D) at (5,0);
    \coordinate (B) at (1.3,2.6);
    \coordinate (C) at (5,2.6);
    \coordinate (N) at (-1.2,0);
    \coordinate (K) at (1.3,1.05);
    \draw[thick] (A) -- (B) -- (C) -- (D) -- cycle;
    \draw[thick] (B) -- (M);
    \draw[thick] (N) -- (C);
    \draw[thick] (N) -- (A);
    \draw[thick] (M)++(0,0.28) -- ++(0.28,0) -- ++(0,-0.28);
    \fill (K) circle (1.6pt);
    \fill (N) circle (1.6pt);
    \node[below] at (A) {$A$};
    \node[below] at (M) {$M$};
    \node[below right] at (D) {$D$};
    \node[above left] at (B) {$B$};
    \node[above right] at (C) {$C$};
    \node[below] at (N) {$N$};
    \node[below right] at (K) {$K$};
\end{tikzpicture}
\end{center}
\vspace{0.2cm}
\noindent
\begin{minipage}[t]{0.42\textwidth}
\textit{Відрізок}\par\vspace{0.2cm}
\textbf{1} \quad $BC$\\[0.25cm]
\textbf{2} \quad $BM$\\[0.25cm]
\textbf{3} \quad $NM$
\end{minipage}
\hfill
\begin{minipage}[t]{0.3\textwidth}
\textit{Довжина відрізка}\par\vspace{0.2cm}
\textbf{А} \quad $8$ см\\[0.15cm]
\textbf{Б} \quad $12$ см\\[0.15cm]
\textbf{В} \quad $15$ см\\[0.15cm]
\textbf{Г} \quad $18$ см\\[0.15cm]
\textbf{Д} \quad $20$ см
\end{minipage}
\hfill
\begin{minipage}[t]{0.2\textwidth}
\vspace{0.4cm}
\begin{flushright}\matchingGrid\end{flushright}
\end{minipage}}
\vspace{0.5cm}
\zadtask{У паралелограмі $ABCD$ бісектриса кута $A = 60^\circ$ перетинає сторону $BC$ у точці $K$, причому $BK : KC = 4 : 3$ і $KC = 6$~см (див. рисунок). Установіть відповідність між відрізком (1--3) та його довжиною (А--Д).
\vspace{0.3cm}
\begin{center}
\begin{tikzpicture}[scale=0.95, thick]
    \coordinate (A) at (0,0);
    \coordinate (B) at (1.4,2.4);
    \coordinate (K) at (4.4,2.4);
    \coordinate (C) at (6.5,2.4);
    \coordinate (D) at (5.1,0);
    \draw[thick] (A) -- (B) -- (C) -- (D) -- cycle;
    \draw[thick] (A) -- (K);
    \fill (K) circle (1.8pt);
    \pic[draw=yearOrange, line width=0.8pt, pic text={\scriptsize $60^\circ$}, angle radius=0.85cm,
         angle eccentricity=1.35] {angle = D--A--K};
    \pic[draw=yearOrange, line width=0.8pt, angle radius=1.15cm] {angle = K--A--B};
    \node[below left] at (A) {$A$};
    \node[above left] at (B) {$B$};
    \node[above] at (K) {$K$};
    \node[above right] at (C) {$C$};
    \node[below right] at (D) {$D$};
\end{tikzpicture}
\end{center}
\vspace{0.2cm}
\noindent
\begin{minipage}[t]{0.42\textwidth}
\textit{Відрізок}\par\vspace{0.15cm}
\textbf{1} \quad $BC$\\[0.25cm]
\textbf{2} \quad $CD$\\[0.25cm]
\textbf{3} \quad $AK$
\end{minipage}
\begin{minipage}[t]{0.30\textwidth}
\textit{Довжина}\par\vspace{0.15cm}
\textbf{А} \quad $8$ см\\[0.1cm]
\textbf{Б} \quad $8\sqrt{3}$ см\\[0.1cm]
\textbf{В} \quad $4\sqrt{3}$ см\\[0.1cm]
\textbf{Г} \quad $12$ см\\[0.1cm]
\textbf{Д} \quad $14$ см
\end{minipage}
\hfill
\begin{minipage}[t]{0.20\textwidth}
\vspace{-0.2cm}
\begin{flushright}\matchingGrid\end{flushright}
\end{minipage}
}
\newpage
\taskBlock{Прямокутник і коло: відповідність відрізків --- сесія 2.06, завдання №18}
\zadtask{Прямокутник $ABCD$ і коло з центром $O$ розташовані так, що коло дотикається до середин бічних сторін $AB$ і $CD$ та перетинає сторону $BC$ у точках $K$ і $M$ (див. рисунок). Відомо, що $BK = 3$ см, $BM = 27$ см. Установіть відповідність між відрізком (1--3) та його довжиною (А--Д).
\vspace{0.3cm}
\begin{center}
\begin{tikzpicture}[scale=0.2, thick]
    \coordinate (A) at (0,0);
    \coordinate (B) at (0,18);
    \coordinate (C) at (30,18);
    \coordinate (D) at (30,0);
    \coordinate (O) at (15,9);
    \coordinate (K) at (3,18);
    \coordinate (M) at (27,18);
    \draw[thick] (A) -- (B) -- (C) -- (D) -- cycle;
    \draw[thick, mainGreen] (O) circle (15);
    \fill (O) circle (4pt); \node[below right] at (O) {$O$};
    \fill (K) circle (4pt); \node[above] at (K) {$K$};
    \fill (M) circle (4pt); \node[above] at (M) {$M$};
    \node[below left] at (A) {$A$};
    \node[above left] at (B) {$B$};
    \node[above right] at (C) {$C$};
    \node[below right] at (D) {$D$};
\end{tikzpicture}
\end{center}
\vspace{0.2cm}
\noindent
\begin{minipage}[t]{0.42\textwidth}
\textit{Відрізок}\par\vspace{0.2cm}
\textbf{1} \quad $KM$\\[0.25cm]
\textbf{2} \quad $KO$\\[0.25cm]
\textbf{3} \quad $AB$
\end{minipage}
\hfill
\begin{minipage}[t]{0.3\textwidth}
\textit{Довжина}\par\vspace{0.2cm}
\textbf{А} \quad $3$ см\\[0.15cm]
\textbf{Б} \quad $15$ см\\[0.15cm]
\textbf{В} \quad $18$ см\\[0.15cm]
\textbf{Г} \quad $24$ см\\[0.15cm]
\textbf{Д} \quad $30$ см
\end{minipage}
\hfill
\begin{minipage}[t]{0.2\textwidth}
\vspace{0.4cm}
\begin{flushright}\matchingGrid\end{flushright}
\end{minipage}}
\hypertarget{task:73}{}\par\vspace{0.05cm}\noindent\hyperlink{sol:73}{\fbox{\footnotesize розв'язок $\rightarrow$}}\par\vspace{0.2cm}
\zadtask{Прямокутник $ABCD$ і коло з центром $O$ розташовані так, що коло дотикається до середин бічних сторін $AB$ і $CD$ та перетинає сторону $BC$ у точках $K$ і $M$ (див. рисунок). Відомо, що $BK = 2$ см, $BM = 50$ см. Установіть відповідність між відрізком (1--3) та його довжиною (А--Д).
\vspace{0.3cm}
\begin{center}
\begin{tikzpicture}[scale=0.12, thick]
    \coordinate (A) at (0,0);
    \coordinate (B) at (0,20);
    \coordinate (C) at (52,20);
    \coordinate (D) at (52,0);
    \coordinate (O) at (26,10);
    \coordinate (K) at (2,20);
    \coordinate (M) at (50,20);
    \draw[thick] (A) -- (B) -- (C) -- (D) -- cycle;
    \draw[thick, mainGreen] (O) circle (26);
    \fill (O) circle (6pt); \node[below right] at (O) {$O$};
    \fill (K) circle (6pt); \node[above] at (K) {$K$};
    \fill (M) circle (6pt); \node[above] at (M) {$M$};
    \node[below left] at (A) {$A$};
    \node[above left] at (B) {$B$};
    \node[above right] at (C) {$C$};
    \node[below right] at (D) {$D$};
\end{tikzpicture}
\end{center}
\vspace{0.2cm}
\noindent
\begin{minipage}[t]{0.42\textwidth}
\textit{Відрізок}\par\vspace{0.2cm}
\textbf{1} \quad $KM$\\[0.25cm]
\textbf{2} \quad $KO$\\[0.25cm]
\textbf{3} \quad $AB$
\end{minipage}
\hfill
\begin{minipage}[t]{0.3\textwidth}
\textit{Довжина}\par\vspace{0.2cm}
\textbf{А} \quad $20$ см\\[0.15cm]
\textbf{Б} \quad $2$ см\\[0.15cm]
\textbf{В} \quad $48$ см\\[0.15cm]
\textbf{Г} \quad $26$ см\\[0.15cm]
\textbf{Д} \quad $52$ см
\end{minipage}
\hfill
\begin{minipage}[t]{0.2\textwidth}
\vspace{0.4cm}
\begin{flushright}\matchingGrid\end{flushright}
\end{minipage}}
\zadtask{На рисунку зображено прямокутник $ABCD$ і коло, яке дотикається до сторони $AB$ та сторін $BC$ і $AD$ у точках $M$ і $K$ відповідно. Периметр чотирикутника $ABMK$ дорівнює $18$~см, а довжина відрізка $KC$~--- $10$~см. Установіть відповідність між величиною (1--3) та її значенням (А--Д).
\vspace{0.3cm}
\begin{center}
\begin{tikzpicture}[scale=0.55, thick]
    \coordinate (A) at (0,0);
    \coordinate (B) at (0,6);
    \coordinate (Cc) at (11,6);
    \coordinate (D) at (11,0);
    \coordinate (M) at (3,6);
    \coordinate (K) at (3,0);
    \coordinate (O) at (3,3);
    \draw[thick] (A) -- (B) -- (Cc) -- (D) -- cycle;
    \draw[thick, mainGreen] (O) circle (3);
    \draw[thick] (M) -- (K);
    \draw[thick] (K) -- (Cc);
    \fill (M) circle (2.2pt); \node[above] at (M) {$M$};
    \fill (K) circle (2.2pt); \node[below] at (K) {$K$};
    \node[below left] at (A) {$A$};
    \node[above left] at (B) {$B$};
    \node[above right] at (Cc) {$C$};
    \node[below right] at (D) {$D$};
\end{tikzpicture}
\end{center}
\vspace{0.2cm}
\noindent
\begin{minipage}[t]{0.46\textwidth}
\textit{Величина}\par\vspace{0.15cm}
\textbf{1} \quad радіус кола\\[0.25cm]
\textbf{2} \quad сторона $BC$\\[0.25cm]
\textbf{3} \quad відстань від точки $M$ до прямої $KC$
\end{minipage}
\begin{minipage}[t]{0.28\textwidth}
\textit{Значення}\par\vspace{0.15cm}
\textbf{А} \quad $3$ см\\[0.1cm]
\textbf{Б} \quad $4{,}8$ см\\[0.1cm]
\textbf{В} \quad $5$ см\\[0.1cm]
\textbf{Г} \quad $8$ см\\[0.1cm]
\textbf{Д} \quad $11$ см
\end{minipage}
\hfill
\begin{minipage}[t]{0.20\textwidth}
\vspace{-0.2cm}
\begin{flushright}\matchingGrid\end{flushright}
\end{minipage}
}
\newpage
\taskBlock{Планіметрія: прямокутник і рівнобедрений трикутник (відповідність) --- сесія 3.06, завдання №18}
\zadnum На рисунку зображено прямокутник $ABCD$ та рівнобедрений трикутник $DKM$ ($DK = KM$). Пряма, що містить діагональ $AC$ прямокутника, проходить через точку $K$. $AB = 16$ см, $AD = 12$ см, $DM = 2AD$. Установіть відповідність між відрізком (1--3) та його довжиною (А--Д).
\vspace{0.3cm}
\begin{center}
\begin{tikzpicture}[scale=0.17, thick]
    \coordinate (A) at (0,0);
    \coordinate (B) at (0,16);
    \coordinate (Cc) at (12,16);
    \coordinate (D) at (12,0);
    \coordinate (M) at (36,0);
    \coordinate (K) at (24,32);
    \draw[thick] (A) -- (B) -- (Cc) -- (D) -- cycle;
    \draw[thick] (A) -- (M);
    \draw[thick] (D) -- (K) -- (M);
    \draw[thick] (A) -- (K);
    \fill (K) circle (3pt); \node[above, font=\scriptsize] at (K) {$K$};
    \node[below left, font=\scriptsize] at (A) {$A$};
    \node[above left, font=\scriptsize] at (B) {$B$};
    \node[above, font=\scriptsize] at (Cc) {$C$};
    \node[below, font=\scriptsize] at (D) {$D$};
    \node[below right, font=\scriptsize] at (M) {$M$};
\end{tikzpicture}
\end{center}
\vspace{0.2cm}
\noindent
\begin{minipage}[t]{0.46\textwidth}
\textit{Початок речення}\par\vspace{0.2cm}
\textbf{1} \quad Довжина відрізка $AM$ дорівнює\\[0.2cm]
\textbf{2} \quad Відстань від точки $K$ до прямої $AM$ дорівнює\\[0.2cm]
\textbf{3} \quad Довжина відрізка $CK$ дорівнює
\end{minipage}
\hfill
\begin{minipage}[t]{0.28\textwidth}
\textit{Закінчення}\par\vspace{0.2cm}
\textbf{А} \quad $20$ см\\[0.15cm]
\textbf{Б} \quad $24$ см\\[0.15cm]
\textbf{В} \quad $32$ см\\[0.15cm]
\textbf{Г} \quad $36$ см\\[0.15cm]
\textbf{Д} \quad $48$ см
\end{minipage}
\hfill
\begin{minipage}[t]{0.2\textwidth}
\vspace{0.5cm}
\begin{flushright}\matchingGrid\end{flushright}
\end{minipage}
\hypertarget{task:76}{}\par\vspace{0.05cm}\noindent\hyperlink{sol:76}{\fbox{\footnotesize розв'язок $\rightarrow$}}\par\vspace{0.2cm}
\vspace{0.7cm}
\zadnum На рисунку зображено прямокутник $ABCD$ та рівнобедрений трикутник $DKM$ ($DK = KM$). Пряма, що містить діагональ $AC$ прямокутника, проходить через точку $K$. $AB = 12$ см, $AD = 9$ см, $DM = 2AD$. Установіть відповідність між відрізком (1--3) та його довжиною (А--Д).
\vspace{0.3cm}
\begin{center}
\begin{tikzpicture}[scale=0.22, thick]
    \coordinate (A) at (0,0);
    \coordinate (B) at (0,12);
    \coordinate (Cc) at (9,12);
    \coordinate (D) at (9,0);
    \coordinate (M) at (27,0);
    \coordinate (K) at (18,24);
    \draw[thick] (A) -- (B) -- (Cc) -- (D) -- cycle;
    \draw[thick] (A) -- (M);
    \draw[thick] (D) -- (K) -- (M);
    \draw[thick] (A) -- (K);
    \fill (K) circle (3pt); \node[above, font=\scriptsize] at (K) {$K$};
    \node[below left, font=\scriptsize] at (A) {$A$};
    \node[above left, font=\scriptsize] at (B) {$B$};
    \node[above, font=\scriptsize] at (Cc) {$C$};
    \node[below, font=\scriptsize] at (D) {$D$};
    \node[below right, font=\scriptsize] at (M) {$M$};
\end{tikzpicture}
\end{center}
\vspace{0.2cm}
\noindent
\begin{minipage}[t]{0.46\textwidth}
\textit{Початок речення}\par\vspace{0.2cm}
\textbf{1} \quad Довжина відрізка $AM$ дорівнює\\[0.2cm]
\textbf{2} \quad Відстань від точки $K$ до прямої $AM$ дорівнює\\[0.2cm]
\textbf{3} \quad Довжина відрізка $CK$ дорівнює
\end{minipage}
\hfill
\begin{minipage}[t]{0.28\textwidth}
\textit{Закінчення}\par\vspace{0.2cm}
\textbf{А} \quad $24$ см\\[0.15cm]
\textbf{Б} \quad $15$ см\\[0.15cm]
\textbf{В} \quad $27$ см\\[0.15cm]
\textbf{Г} \quad $30$ см\\[0.15cm]
\textbf{Д} \quad $36$ см
\end{minipage}
\hfill
\begin{minipage}[t]{0.2\textwidth}
\vspace{0.5cm}
\begin{flushright}\matchingGrid\end{flushright}
\end{minipage}
\vspace{0.7cm}
\zadnum У ромб $ABCD$ вписано квадрат $KLMN$, сторона $KL$ якого перетинає діагональ $AC$ у точці $P$ (див. рисунок). $AL = 10$~см, $AP = 8$~см. Установіть відповідність між величиною (1--3) та її значенням (А--Д).
\vspace{0.3cm}
\begin{center}
\begin{tikzpicture}[scale=0.45, thick]
    \coordinate (A) at (-7,0);
    \coordinate (B) at (0,5.25);
    \coordinate (C) at (7,0);
    \coordinate (D) at (0,-5.25);
    \coordinate (K) at (-3,-3);
    \coordinate (L) at (-3,3);
    \coordinate (M) at (3,3);
    \coordinate (N) at (3,-3);
    \coordinate (P) at (-3,0);
    \draw[thick] (A) -- (B) -- (C) -- (D) -- cycle;
    \draw[thick, mainGreen] (K) -- (L) -- (M) -- (N) -- cycle;
    \draw[thick, dashed] (A) -- (C);
    \fill (P) circle (2.2pt);
    \node[below left] at (A) {$A$};
    \node[above] at (B) {$B$};
    \node[below right] at (C) {$C$};
    \node[below] at (D) {$D$};
    \node[below left] at (K) {$K$};
    \node[above left] at (L) {$L$};
    \node[above right] at (M) {$M$};
    \node[below right] at (N) {$N$};
    \node[above right] at (P) {$P$};
\end{tikzpicture}
\end{center}
\vspace{0.2cm}
\noindent
\begin{minipage}[t]{0.46\textwidth}
\textit{Величина}\par\vspace{0.15cm}
\textbf{1} \quad сторона квадрата $KLMN$\\[0.25cm]
\textbf{2} \quad діагональ $BD$ ромба\\[0.25cm]
\textbf{3} \quad радіус кола, вписаного в ромб
\end{minipage}
\begin{minipage}[t]{0.28\textwidth}
\textit{Значення}\par\vspace{0.15cm}
\textbf{А} \quad $8{,}4$ см\\[0.1cm]
\textbf{Б} \quad $12$ см\\[0.1cm]
\textbf{В} \quad $6$ см\\[0.1cm]
\textbf{Г} \quad $21$ см\\[0.1cm]
\textbf{Д} \quad $10$ см
\end{minipage}
\hfill
\begin{minipage}[t]{0.20\textwidth}
\vspace{0.5cm}
\begin{flushright}\matchingGrid\end{flushright}
\end{minipage}
\newpage
\taskBlock{Рівні прямокутні рівнобедрені трикутники (відповідність) --- сесія 4.06, завдання №18}
\zadtask{Прямокутні рівнобедрені трикутники $ABC$ і $KMN$ є рівними (див. рисунок). Точки $K$, $C$, $N$ і $A$ лежать на прямій $AK$, причому $AN = NC = CK$. $AC = 12$~см. Установіть відповідність між відрізком (1--3) та його довжиною (А--Д).
\vspace{0.2cm}
\noindent
\begin{minipage}[c]{0.62\textwidth}
\begin{minipage}[t]{0.55\textwidth}
\textit{Відрізок}\par\vspace{0.2cm}
\textbf{1} \quad $AK$\\[0.25cm]
\textbf{2} \quad $AB$\\[0.25cm]
\textbf{3} \quad $BM$
\end{minipage}%
\hfill
\begin{minipage}[t]{0.42\textwidth}
\textit{Довжина відрізка}\par\vspace{0.2cm}
\textbf{А} \quad $6\sqrt{2}$~см\\[0.15cm]
\textbf{Б} \quad $6\sqrt{5}$~см\\[0.15cm]
\textbf{В} \quad $12\sqrt{2}$~см\\[0.15cm]
\textbf{Г} \quad $18$~см\\[0.15cm]
\textbf{Д} \quad $24$~см
\end{minipage}
\par\vspace{0.3cm}
\matchingGrid
\end{minipage}%
\hfill
\begin{minipage}[c]{0.36\textwidth}
\centering
\begin{tikzpicture}[scale=0.7, thick]
    \coordinate (K) at (-3,3);
    \coordinate (M) at (1,3);
    \coordinate (N) at (1,-1);
    \coordinate (C) at (-1,1);
    \coordinate (B) at (-1,-3);
    \coordinate (A) at (3,-3);
    \draw (K) -- (M) -- (N);
    \draw (C) -- (B) -- (A);
    \draw (K) -- (A);
    \draw (M) ++(-0.4,0) -- ++(0,-0.4) -- ++(0.4,0);
    \draw (B) ++(0,0.4) -- ++(0.4,0) -- ++(0,-0.4);
    \draw (-2-0.15, 2-0.15) -- (-2+0.15, 2+0.15);
    \draw (0-0.15, 0-0.15) -- (0+0.15, 0+0.15);
    \draw (2-0.15, -2-0.15) -- (2+0.15, -2+0.15);
    \fill (C) circle (2pt);
    \fill (N) circle (2pt);
    \node[above left] at (K) {$K$};
    \node[above right] at (M) {$M$};
    \node[left] at (-1.1, 0.9) {$C$};
    \node[right] at (1.1, -0.9) {$N$};
    \node[below left] at (B) {$B$};
    \node[below right] at (A) {$A$};
\end{tikzpicture}
\end{minipage}
}
\hypertarget{task:79}{}\par\vspace{0.05cm}\noindent\hyperlink{sol:79}{\fbox{\footnotesize розв'язок $\rightarrow$}}\par\vspace{0.2cm}
\zadtask{Прямокутні рівнобедрені трикутники $ABC$ і $KMN$ є рівними (див. рисунок). Точки $K$, $C$, $N$ і $A$ лежать на прямій $AK$, причому $AN = NC = CK$. $AC = 8$~см. Установіть відповідність між відрізком (1--3) та його довжиною (А--Д).
\vspace{0.2cm}
\noindent
\begin{minipage}[c]{0.62\textwidth}
\begin{minipage}[t]{0.55\textwidth}
\textit{Відрізок}\par\vspace{0.2cm}
\textbf{1} \quad $AK$\\[0.25cm]
\textbf{2} \quad $AB$\\[0.25cm]
\textbf{3} \quad $BM$
\end{minipage}%
\hfill
\begin{minipage}[t]{0.42\textwidth}
\textit{Довжина відрізка}\par\vspace{0.2cm}
\textbf{А} \quad $4\sqrt{2}$~см\\[0.15cm]
\textbf{Б} \quad $4\sqrt{5}$~см\\[0.15cm]
\textbf{В} \quad $8\sqrt{2}$~см\\[0.15cm]
\textbf{Г} \quad $12$~см\\[0.15cm]
\textbf{Д} \quad $16$~см
\end{minipage}
\par\vspace{0.3cm}
\matchingGrid
\end{minipage}%
\hfill
\begin{minipage}[c]{0.36\textwidth}
\centering
\begin{tikzpicture}[scale=0.7, thick]
    \coordinate (K) at (-3,3);
    \coordinate (M) at (1,3);
    \coordinate (N) at (1,-1);
    \coordinate (C) at (-1,1);
    \coordinate (B) at (-1,-3);
    \coordinate (A) at (3,-3);
    \draw (K) -- (M) -- (N);
    \draw (C) -- (B) -- (A);
    \draw (K) -- (A);
    \draw (M) ++(-0.4,0) -- ++(0,-0.4) -- ++(0.4,0);
    \draw (B) ++(0,0.4) -- ++(0.4,0) -- ++(0,-0.4);
    \draw (-2-0.15, 2-0.15) -- (-2+0.15, 2+0.15);
    \draw (0-0.15, 0-0.15) -- (0+0.15, 0+0.15);
    \draw (2-0.15, -2-0.15) -- (2+0.15, -2+0.15);
    \fill (C) circle (2pt);
    \fill (N) circle (2pt);
    \node[above left] at (K) {$K$};
    \node[above right] at (M) {$M$};
    \node[left] at (-1.1, 0.9) {$C$};
    \node[right] at (1.1, -0.9) {$N$};
    \node[below left] at (B) {$B$};
    \node[below right] at (A) {$A$};
\end{tikzpicture}
\end{minipage}
}
\zadtask{На рисунку зображено два прямокутники $ABCD$ та $CEFG$, розташовані так, що точки $D$, $C$, $E$ лежать на одній прямій, а точки $B$, $C$, $G$~--- на іншій. Відношення довжини до ширини в кожному прямокутнику дорівнює $4 : 3$. Периметр $ABCD$ дорівнює $28$~см, а периметр $CEFG$~--- $42$~см. Установіть відповідність між відрізком (1--3) та його довжиною (А--Д).
\vspace{0.3cm}
\begin{center}
\begin{tikzpicture}[scale=0.3, thick]
    \draw (0,0) rectangle (8,6);
    \draw (8,6) rectangle (20,15);
    \node[below left] at (0,0) {$A$};
    \node[below right] at (8,0) {$B$};
    \node[above left] at (0,6) {$D$};
    \node[below right] at (8,6) {$C$};
    \node[below right] at (20,6) {$E$};
    \node[above right] at (20,15) {$F$};
    \node[above left] at (8,15) {$G$};
    \foreach \p in {(0,0),(8,0),(8,6),(0,6),(20,6),(20,15),(8,15)} {\fill \p circle (5pt);}
\end{tikzpicture}
\end{center}
\vspace{0.2cm}
\noindent
\begin{minipage}[t]{0.42\textwidth}
\textit{Відрізок}\par\vspace{0.15cm}
\textbf{1} \quad діагональ $AC$\\[0.25cm]
\textbf{2} \quad відрізок $DE$\\[0.25cm]
\textbf{3} \quad відрізок $AF$
\end{minipage}
\begin{minipage}[t]{0.30\textwidth}
\textit{Довжина відрізка}\par\vspace{0.15cm}
\textbf{А} \quad $10$~см\\[0.1cm]
\textbf{Б} \quad $15$~см\\[0.1cm]
\textbf{В} \quad $17$~см\\[0.1cm]
\textbf{Г} \quad $20$~см\\[0.1cm]
\textbf{Д} \quad $25$~см
\end{minipage}
\hfill
\begin{minipage}[t]{0.20\textwidth}
\vspace{-0.2cm}
\begin{flushright}\matchingGrid\end{flushright}
\end{minipage}
}
\newpage
\taskBlock{Кути в комбінації фігур (відповідність) --- сесія 5.06, завдання №18}
\zadnum На рисунку зображено правильний трикутник $ABC$ та прямокутні рівнобедрені трикутники $AKB$ і $BMC$. Установіть відповідність між кутом (1--3) та його градусною мірою (А--Д).
\vspace{0.3cm}
\begin{center}
\begin{tikzpicture}[scale=1.9, thick]
    \coordinate (A) at (0,0);
    \coordinate (C) at (2,0);
    \coordinate (B) at (1,1.732);
    \coordinate (K) at (-0.366,1.366);
    \coordinate (M) at (2.366,1.366);
    \draw[thick] (A) -- (B) -- (C) -- cycle;
    \draw[thick] (A) -- (K) -- (B);
    \draw[thick] (B) -- (M) -- (C);
    \draw[thick] (A) -- (M);
    \draw[thick] (K) -- (B);
    \pic[draw=yearOrange, line width=0.5pt, angle radius=0.3cm] {right angle = A--K--B};
    \pic[draw=yearOrange, line width=0.5pt, angle radius=0.3cm] {right angle = B--M--C};
    \node[below left, font=\scriptsize] at (A) {$A$};
    \node[above, font=\scriptsize] at (B) {$B$};
    \node[below right, font=\scriptsize] at (C) {$C$};
    \node[left, font=\scriptsize] at (K) {$K$};
    \node[right, font=\scriptsize] at (M) {$M$};
\end{tikzpicture}
\end{center}
\vspace{0.2cm}
\noindent
\begin{minipage}[t]{0.32\textwidth}
\textit{Кут}\par\vspace{0.2cm}
\textbf{1} \quad $\angle ACM$\\[0.25cm]
\textbf{2} \quad $\angle KBM$\\[0.25cm]
\textbf{3} \quad $\angle KAM$
\end{minipage}
\hfill
\begin{minipage}[t]{0.34\textwidth}
\textit{Градусна міра}\par\vspace{0.2cm}
\textbf{А} \quad $75^\circ$\\[0.15cm]
\textbf{Б} \quad $90^\circ$\\[0.15cm]
\textbf{В} \quad $105^\circ$\\[0.15cm]
\textbf{Г} \quad $135^\circ$\\[0.15cm]
\textbf{Д} \quad $150^\circ$
\end{minipage}
\hfill
\begin{minipage}[t]{0.2\textwidth}
\vspace{0.5cm}
\begin{flushright}\matchingGrid\end{flushright}
\end{minipage}
\par\vspace{0.5cm}
\hypertarget{task:82}{}\par\vspace{0.05cm}\noindent\hyperlink{sol:82}{\fbox{\footnotesize розв'язок $\rightarrow$}}\par\vspace{0.2cm}
\zadnum На рисунку зображено правильний трикутник $ABC$ та прямокутні рівнобедрені трикутники $AKB$ і $BMC$. Установіть відповідність між кутом (1--3) та його градусною мірою (А--Д).
\vspace{0.3cm}
\begin{center}
\begin{tikzpicture}[scale=1.9, thick]
    \coordinate (A) at (0,0);
    \coordinate (C) at (2,0);
    \coordinate (B) at (1,1.732);
    \coordinate (K) at (-0.366,1.366);
    \coordinate (M) at (2.366,1.366);
    \draw[thick] (A) -- (B) -- (C) -- cycle;
    \draw[thick] (A) -- (K) -- (B);
    \draw[thick] (B) -- (M) -- (C);
    \draw[thick] (K) -- (M);
    \draw[thick] (A) -- (K);
    \pic[draw=yearOrange, line width=0.5pt, angle radius=0.3cm] {right angle = A--K--B};
    \pic[draw=yearOrange, line width=0.5pt, angle radius=0.3cm] {right angle = B--M--C};
    \node[below left, font=\scriptsize] at (A) {$A$};
    \node[above, font=\scriptsize] at (B) {$B$};
    \node[below right, font=\scriptsize] at (C) {$C$};
    \node[left, font=\scriptsize] at (K) {$K$};
    \node[right, font=\scriptsize] at (M) {$M$};
\end{tikzpicture}
\end{center}
\vspace{0.2cm}
\noindent
\begin{minipage}[t]{0.32\textwidth}
\textit{Кут}\par\vspace{0.2cm}
\textbf{1} \quad $\angle MBC$\\[0.25cm]
\textbf{2} \quad $\angle AKM$\\[0.25cm]
\textbf{3} \quad $\angle KAC$
\end{minipage}
\hfill
\begin{minipage}[t]{0.34\textwidth}
\textit{Градусна міра}\par\vspace{0.2cm}
\textbf{А} \quad $45^\circ$\\[0.15cm]
\textbf{Б} \quad $75^\circ$\\[0.15cm]
\textbf{В} \quad $105^\circ$\\[0.15cm]
\textbf{Г} \quad $120^\circ$\\[0.15cm]
\textbf{Д} \quad $135^\circ$
\end{minipage}
\hfill
\begin{minipage}[t]{0.2\textwidth}
\vspace{0.5cm}
\begin{flushright}\matchingGrid\end{flushright}
\end{minipage}
\par\vspace{0.5cm}
\zadnum На рисунку зображено квадрат $ABCD$ та правильний трикутник $ABK$, побудований \textit{усередині} квадрата на стороні $AB$. Установіть відповідність між кутом (1--3) та його градусною мірою (А--Д).
\vspace{0.3cm}
\begin{center}
\begin{tikzpicture}[scale=1.9, thick]
    \coordinate (A) at (0,0);
    \coordinate (B) at (2,0);
    \coordinate (C) at (2,-2);
    \coordinate (D) at (0,-2);
    \coordinate (K) at (1,-1.732);
    \draw[thick] (A) -- (B) -- (C) -- (D) -- cycle;
    \draw[thick] (A) -- (K) -- (B);
    \draw[thick] (K) -- (C);
    \draw[thick] (K) -- (D);
    \node[above left, font=\scriptsize] at (A) {$A$};
    \node[above right, font=\scriptsize] at (B) {$B$};
    \node[below right, font=\scriptsize] at (C) {$C$};
    \node[below left, font=\scriptsize] at (D) {$D$};
    \node[below, font=\scriptsize] at (K) {$K$};
\end{tikzpicture}
\end{center}
\vspace{0.2cm}
\noindent
\begin{minipage}[t]{0.32\textwidth}
\textit{Кут}\par\vspace{0.2cm}
\textbf{1} \quad $\angle KCB$\\[0.25cm]
\textbf{2} \quad $\angle KCD$\\[0.25cm]
\textbf{3} \quad $\angle DKC$
\end{minipage}
\hfill
\begin{minipage}[t]{0.34\textwidth}
\textit{Градусна міра}\par\vspace{0.2cm}
\textbf{А} \quad $15^\circ$\\[0.15cm]
\textbf{Б} \quad $75^\circ$\\[0.15cm]
\textbf{В} \quad $105^\circ$\\[0.15cm]
\textbf{Г} \quad $135^\circ$\\[0.15cm]
\textbf{Д} \quad $150^\circ$
\end{minipage}
\hfill
\begin{minipage}[t]{0.2\textwidth}
\vspace{0.5cm}
\begin{flushright}\matchingGrid\end{flushright}
\end{minipage}
\par\vspace{0.5cm}
\newpage
\chapterTitle{РОЗВ\textquotesingle ЯЗКИ ОРИГІНАЛЬНИХ ЗАДАЧ}
\noindent\small У цьому розділі зібрано стислі розв\textquotesingle язки до перших задач кожного блоку (оригіналів НМТ-2026). Натискайте посилання \fbox{\footnotesize $\leftarrow$ Задача №$N$}, щоб повернутися до задачі.
\par\vspace{0.4cm}
\par\vspace{0.3cm}\noindent\hypertarget{sol:1}{}
\begin{tcolorbox}[colback=titleBg!40, colframe=mainGreen!55, boxrule=0.4pt, arc=2pt, left=6pt, right=6pt, top=4pt, bottom=4pt, breakable]
\noindent\small\textbf{Задача №1.}\quad Сума кутів трикутника $180^\circ$: $\angle C=180^\circ-30^\circ-35^\circ=115^\circ$. Відповідь: \textbf{В}.
\par\vspace{0.1cm}\hfill\hyperlink{task:1}{\fbox{\footnotesize $\leftarrow$ Задача №1}}
\end{tcolorbox}
\par\vspace{0.3cm}\noindent\hypertarget{sol:4}{}
\begin{tcolorbox}[colback=titleBg!40, colframe=mainGreen!55, boxrule=0.4pt, arc=2pt, left=6pt, right=6pt, top=4pt, bottom=4pt, breakable]
\noindent\small\textbf{Задача №4.}\quad Гострий кут ромба $180^\circ-140^\circ=40^\circ$. Більша діагональ сполучає вершини тупих кутів і ділить навпіл гострий кут, тому з стороною утворює $40^\circ:2=20^\circ$. Відповідь: \textbf{А} ($20^\circ$).
\par\vspace{0.1cm}\hfill\hyperlink{task:4}{\fbox{\footnotesize $\leftarrow$ Задача №4}}
\end{tcolorbox}
\par\vspace{0.3cm}\noindent\hypertarget{sol:7}{}
\begin{tcolorbox}[colback=titleBg!40, colframe=mainGreen!55, boxrule=0.4pt, arc=2pt, left=6pt, right=6pt, top=4pt, bottom=4pt, breakable]
\noindent\small\textbf{Задача №7.}\quad Розгорнутий кут півкола $180^\circ$ поділено на $5$ рівних частин, тож кожна частина (і кут $\alpha$ між сусідніми радіусами) дорівнює $\dfrac{180^\circ}{5}=36^\circ$. Відповідь: \textbf{В}.
\par\vspace{0.1cm}\hfill\hyperlink{task:7}{\fbox{\footnotesize $\leftarrow$ Задача №7}}
\end{tcolorbox}
\par\vspace{0.3cm}\noindent\hypertarget{sol:10}{}
\begin{tcolorbox}[colback=titleBg!40, colframe=mainGreen!55, boxrule=0.4pt, arc=2pt, left=6pt, right=6pt, top=4pt, bottom=4pt, breakable]
\noindent\small\textbf{Задача №10.}\quad Оскільки $AB=BC$, кути при основі рівні: $\angle A=\angle C$. За умовою $\angle B=2\angle C$. Сума кутів: $\angle C+\angle C+2\angle C=180^\circ$, $4\angle C=180^\circ$, $\angle C=45^\circ$, тож $\angle B=90^\circ$. Відповідь: \textbf{Г}.
\par\vspace{0.1cm}\hfill\hyperlink{task:10}{\fbox{\footnotesize $\leftarrow$ Задача №10}}
\end{tcolorbox}
\par\vspace{0.3cm}\noindent\hypertarget{sol:13}{}
\begin{tcolorbox}[colback=titleBg!40, colframe=mainGreen!55, boxrule=0.4pt, arc=2pt, left=6pt, right=6pt, top=4pt, bottom=4pt, breakable]
\noindent\small\textbf{Задача №13.}\quad Точки $A$, $B$, $K$ лежать на одній прямій, тому $\angle DKA = 180^\circ - \angle BKD = 180^\circ - 128^\circ = 52^\circ$. У трикутнику $AKD$: $\angle BAD = 180^\circ - \angle ADK - \angle DKA = 180^\circ - 68^\circ - 52^\circ = 60^\circ$. Відповідь: \textbf{Г} ($60^\circ$).
\par\vspace{0.1cm}\hfill\hyperlink{task:13}{\fbox{\footnotesize $\leftarrow$ Задача №13}}
\end{tcolorbox}
\par\vspace{0.3cm}\noindent\hypertarget{sol:16}{}
\begin{tcolorbox}[colback=titleBg!40, colframe=mainGreen!55, boxrule=0.4pt, arc=2pt, left=6pt, right=6pt, top=4pt, bottom=4pt, breakable]
\noindent\small\textbf{Задача №16.}\quad Нехай $BK=x$, тоді $AK=x+5$, а $AK+BK=AB$: $x+5+x=21$, $2x=16$, $x=8$. Отже $AK=8+5=13$~см. Відповідь: \textbf{Д}.
\par\vspace{0.1cm}\hfill\hyperlink{task:16}{\fbox{\footnotesize $\leftarrow$ Задача №16}}
\end{tcolorbox}
\par\vspace{0.3cm}\noindent\hypertarget{sol:19}{}
\begin{tcolorbox}[colback=titleBg!40, colframe=mainGreen!55, boxrule=0.4pt, arc=2pt, left=6pt, right=6pt, top=4pt, bottom=4pt, breakable]
\noindent\small\textbf{Задача №19.}\quad $\angle BKA=180^\circ-120^\circ=60^\circ$. У прямокутному трикутнику $ABK$ ($\angle A=90^\circ$): $AB=BK\sin 60^\circ=12\cdot\dfrac{\sqrt3}{2}=6\sqrt3$ --- сторона квадрата. Радіус описаного кола $R=\dfrac{a\sqrt2}{2}=\dfrac{6\sqrt3\cdot\sqrt2}{2}=3\sqrt6$ см. Відповідь: \textbf{Д}.
\par\vspace{0.1cm}\hfill\hyperlink{task:19}{\fbox{\footnotesize $\leftarrow$ Задача №19}}
\end{tcolorbox}
\par\vspace{0.3cm}\noindent\hypertarget{sol:22}{}
\begin{tcolorbox}[colback=titleBg!40, colframe=mainGreen!55, boxrule=0.4pt, arc=2pt, left=6pt, right=6pt, top=4pt, bottom=4pt, breakable]
\noindent\small\textbf{Задача №22.}\quad Висота трапеції $h=AB\sin60^\circ=6\cdot\dfrac{\sqrt3}{2}=3\sqrt3$. У прямокутному $\triangle ABP$ ($BP\perp AB$, $\angle A=60^\circ$) маємо $AP=\dfrac{AB}{\cos60^\circ}=\dfrac{6}{0{,}5}=12$, тож більша основа $AD=2\cdot AP=24$. Проєкція бічної сторони $AB\cos60^\circ=3$, тому менша основа $BC=AD-2\cdot3=18$. Площа $S=\dfrac{AD+BC}{2}\cdot h=\dfrac{24+18}{2}\cdot3\sqrt3=63\sqrt3$. Відповідь: \textbf{Б}.
\par\vspace{0.1cm}\hfill\hyperlink{task:22}{\fbox{\footnotesize $\leftarrow$ Задача №22}}
\end{tcolorbox}
\par\vspace{0.3cm}\noindent\hypertarget{sol:25}{}
\begin{tcolorbox}[colback=titleBg!40, colframe=mainGreen!55, boxrule=0.4pt, arc=2pt, left=6pt, right=6pt, top=4pt, bottom=4pt, breakable]
\noindent\small\textbf{Задача №25.}\quad Радіус $r$ дорівнює відстані від $O$ до $AD$: $r = OD\sin30^\circ = 24\cdot\dfrac12 = 12$ см, а $OK = r = 12$. Сторона $CD$ (висота прямокутника) дорівнює діаметру: $CD = 2r = 24$ см. Відстань між паралельними $OK$ і $CD$ дорівнює $OD\cos30^\circ = 24\cdot\dfrac{\sqrt3}{2} = 12\sqrt3$. $OKCD$~--- прямокутна трапеція: $S = \dfrac{OK + CD}{2}\cdot 12\sqrt3 = \dfrac{12+24}{2}\cdot 12\sqrt3 = 216\sqrt3$ см$^2$. Відповідь: \textbf{В}.
\par\vspace{0.1cm}\hfill\hyperlink{task:25}{\fbox{\footnotesize $\leftarrow$ Задача №25}}
\end{tcolorbox}
\par\vspace{0.3cm}\noindent\hypertarget{sol:28}{}
\begin{tcolorbox}[colback=titleBg!40, colframe=mainGreen!55, boxrule=0.4pt, arc=2pt, left=6pt, right=6pt, top=4pt, bottom=4pt, breakable]
\noindent\small\textbf{Задача №28.}\quad У прямокутному $\triangle DKM$: $KM=DK\sin30^\circ=12\cdot\dfrac{1}{2}=6$, тож висота $BM=2KM=12$, а $DM=DK\cos30^\circ=12\cdot\dfrac{\sqrt3}{2}=6\sqrt3$. У прямокутному $\triangle ABM$: $AM=\sqrt{AB^2-BM^2}=\sqrt{192-144}=4\sqrt3$. Тоді основа $AD=AM+MD=10\sqrt3$, а площа $S=AD\cdot BM=10\sqrt3\cdot12=120\sqrt3$~см$^2$. Відповідь: \textbf{А}.
\par\vspace{0.1cm}\hfill\hyperlink{task:28}{\fbox{\footnotesize $\leftarrow$ Задача №28}}
\end{tcolorbox}
\par\vspace{0.3cm}\noindent\hypertarget{sol:31}{}
\begin{tcolorbox}[colback=titleBg!40, colframe=mainGreen!55, boxrule=0.4pt, arc=2pt, left=6pt, right=6pt, top=4pt, bottom=4pt, breakable]
\noindent\small\textbf{Задача №31.}\quad У прямокутному $\triangle ADK$ ($\angle ADK=90^\circ$): $DK=AK\sin30^\circ=6$, $AD=AK\cos30^\circ=6\sqrt3$. Точки $D,K,C$ лежать на правій стороні, тож $CK=CD-DK=AB-DK=15-6=9$. Кути при $K$: $\angle AKD=60^\circ$, $\angle AKM=90^\circ$, тому $\angle MKC=180^\circ-60^\circ-90^\circ=30^\circ$. У прямокутному $\triangle KCM$ ($\angle C=90^\circ$): $KM=\dfrac{CK}{\cos30^\circ}=\dfrac{9}{\sqrt3/2}=6\sqrt3$. Площа (прямий кут при $K$): $S=\tfrac12\cdot AK\cdot KM=\tfrac12\cdot12\cdot6\sqrt3=36\sqrt3$ см$^2$. Відповідь: \textbf{Д}.
\par\vspace{0.1cm}\hfill\hyperlink{task:31}{\fbox{\footnotesize $\leftarrow$ Задача №31}}
\end{tcolorbox}
\par\vspace{0.3cm}\noindent\hypertarget{sol:34}{}
\begin{tcolorbox}[colback=titleBg!40, colframe=mainGreen!55, boxrule=0.4pt, arc=2pt, left=6pt, right=6pt, top=4pt, bottom=4pt, breakable]
\noindent\small\textbf{Задача №34.}\quad Площа квадрата $6\cdot 6=36$ см$^2$, площа прямокутника $4\cdot 9=36$ см$^2$. Спільний рівнобедрений трикутник має основу, рівну стороні прямокутника $4$ см, і висоту $2$ см, тож його площа $\dfrac{1}{2}\cdot 4\cdot 2=4$ см$^2$. Площа трафарету $=36+36-4=68$ см$^2$. Відповідь: \textbf{А}.
\par\vspace{0.1cm}\hfill\hyperlink{task:34}{\fbox{\footnotesize $\leftarrow$ Задача №34}}
\end{tcolorbox}
\par\vspace{0.3cm}\noindent\hypertarget{sol:37}{}
\begin{tcolorbox}[colback=titleBg!40, colframe=mainGreen!55, boxrule=0.4pt, arc=2pt, left=6pt, right=6pt, top=4pt, bottom=4pt, breakable]
\noindent\small\textbf{Задача №37.}\quad Чотири вирізані чверті кола радіусом $5$ разом дають одне коло, а два долучені півкола радіусом $5$ — також одне коло; їх площі взаємно компенсуються. Тому площа фігури дорівнює площі прямокутника $ABCD$ зі сторонами $50$~м і $20$~м: $50\cdot20=1000$ (м$^2$). Відповідь: \textbf{В} ($1000$).
\par\vspace{0.1cm}\hfill\hyperlink{task:37}{\fbox{\footnotesize $\leftarrow$ Задача №37}}
\end{tcolorbox}
\par\vspace{0.3cm}\noindent\hypertarget{sol:40}{}
\begin{tcolorbox}[colback=titleBg!40, colframe=mainGreen!55, boxrule=0.4pt, arc=2pt, left=6pt, right=6pt, top=4pt, bottom=4pt, breakable]
\noindent\small\textbf{Задача №40.}\quad Площа кільця $=$ площа більшого круга мінус площа меншого: $S=\pi R^2-\pi\cdot3^2=\pi R^2-9\pi$. Відповідь: \textbf{Д}.
\par\vspace{0.1cm}\hfill\hyperlink{task:40}{\fbox{\footnotesize $\leftarrow$ Задача №40}}
\end{tcolorbox}
\par\vspace{0.3cm}\noindent\hypertarget{sol:43}{}
\begin{tcolorbox}[colback=titleBg!40, colframe=mainGreen!55, boxrule=0.4pt, arc=2pt, left=6pt, right=6pt, top=4pt, bottom=4pt, breakable]
\noindent\small\textbf{Задача №43.}\quad I~--- істинне (січна через центр перетинає коло у двох точках). II~--- істинне (діаметр, перпендикулярний до хорди, ділить її навпіл). III~--- хибне (будь-які два діаметри перетинаються в центрі). Істинні I і II. Відповідь: \textbf{В}.
\par\vspace{0.1cm}\hfill\hyperlink{task:43}{\fbox{\footnotesize $\leftarrow$ Задача №43}}
\end{tcolorbox}
\par\vspace{0.3cm}\noindent\hypertarget{sol:46}{}
\begin{tcolorbox}[colback=titleBg!40, colframe=mainGreen!55, boxrule=0.4pt, arc=2pt, left=6pt, right=6pt, top=4pt, bottom=4pt, breakable]
\noindent\small\textbf{Задача №46.}\quad \textbf{I}: найбільший кут прямокутного трикутника дорівнює $90^\circ$, а половина розгорнутого кута — $180^\circ:2=90^\circ$ (правильно). \textbf{II}: у рівносторонньому трикутнику всі кути по $60^\circ$, сума двох — $120^\circ$ (правильно). \textbf{III}: гострокутний трикутник може мати кут $\le60^\circ$ (напр. $50^\circ,60^\circ,70^\circ$), тож неправильно. Відповідь: \textbf{В} (лише I та II).
\par\vspace{0.1cm}\hfill\hyperlink{task:46}{\fbox{\footnotesize $\leftarrow$ Задача №46}}
\end{tcolorbox}
\par\vspace{0.3cm}\noindent\hypertarget{sol:49}{}
\begin{tcolorbox}[colback=titleBg!40, colframe=mainGreen!55, boxrule=0.4pt, arc=2pt, left=6pt, right=6pt, top=4pt, bottom=4pt, breakable]
\noindent\small\textbf{Задача №49.}\quad У рівнобедреному трикутнику бісектриса, проведена до основи, є також висотою й медіаною, тому $BK\perp AC$ (II) і $AK=KC$ (III). Рівність $AB+BC=AC$ (I) хибна (нерівність трикутника). Правильні II та III. Відповідь: \textbf{Г}.
\par\vspace{0.1cm}\hfill\hyperlink{task:49}{\fbox{\footnotesize $\leftarrow$ Задача №49}}
\end{tcolorbox}
\par\vspace{0.3cm}\noindent\hypertarget{sol:52}{}
\begin{tcolorbox}[colback=titleBg!40, colframe=mainGreen!55, boxrule=0.4pt, arc=2pt, left=6pt, right=6pt, top=4pt, bottom=4pt, breakable]
\noindent\small\textbf{Задача №52.}\quad I -- правильно (у правильному трикутнику всі медіани рівні). II -- хибне: медіана ділить лише сторону навпіл, периметри частин зазвичай різні. III -- правильно: медіана ділить трикутник на два рівновеликі (рівні за площею). Отже правильні I та III. Відповідь: \textbf{В}.
\par\vspace{0.1cm}\hfill\hyperlink{task:52}{\fbox{\footnotesize $\leftarrow$ Задача №52}}
\end{tcolorbox}
\par\vspace{0.3cm}\noindent\hypertarget{sol:55}{}
\begin{tcolorbox}[colback=titleBg!40, colframe=mainGreen!55, boxrule=0.4pt, arc=2pt, left=6pt, right=6pt, top=4pt, bottom=4pt, breakable]
\noindent\small\textbf{Задача №55.}\quad I~--- хибне: хорда може бути меншою за радіус. II~--- правильне: кінці діаметра ділять коло на дві рівні півкола. III~--- правильне: рівні хорди стягують рівні дуги. Отже правильні лише II та III. Відповідь: \textbf{Г}.
\par\vspace{0.1cm}\hfill\hyperlink{task:55}{\fbox{\footnotesize $\leftarrow$ Задача №55}}
\end{tcolorbox}
\par\vspace{0.3cm}\noindent\hypertarget{sol:58}{}
\begin{tcolorbox}[colback=titleBg!40, colframe=mainGreen!55, boxrule=0.4pt, arc=2pt, left=6pt, right=6pt, top=4pt, bottom=4pt, breakable]
\noindent\small\textbf{Задача №58.}\quad I~--- правильне (точка дотику вписаного кола з основою рівнобедреного трикутника збігається з її серединою). II~--- хибне: центр описаного кола тупокутного трикутника лежить \emph{поза} трикутником. III~--- правильне (вписаний кут, що спирається на діаметр, прямий). Правильні I та III. Відповідь: \textbf{В}.
\par\vspace{0.1cm}\hfill\hyperlink{task:58}{\fbox{\footnotesize $\leftarrow$ Задача №58}}
\end{tcolorbox}
\par\vspace{0.3cm}\noindent\hypertarget{sol:61}{}
\begin{tcolorbox}[colback=titleBg!40, colframe=mainGreen!55, boxrule=0.4pt, arc=2pt, left=6pt, right=6pt, top=4pt, bottom=4pt, breakable]
\noindent\small\textbf{Задача №61.}\quad $MN$~--- середня лінія, тож $MN\parallel AC$ і $MN=\tfrac12 AC$: $AMNC$~--- трапеція (I правильне). Тричикутник $MBN$ рівнобедрений лише за умови $AB=BC$, тож для довільного трикутника твердження II хибне. Периметр $MBN=\tfrac12(AB+BC+AC)=\tfrac12 P_{ABC}$ (III правильне). Відповідь: \textbf{В}.
\par\vspace{0.1cm}\hfill\hyperlink{task:61}{\fbox{\footnotesize $\leftarrow$ Задача №61}}
\end{tcolorbox}
\par\vspace{0.3cm}\noindent\hypertarget{sol:64}{}
\begin{tcolorbox}[colback=titleBg!40, colframe=mainGreen!55, boxrule=0.4pt, arc=2pt, left=6pt, right=6pt, top=4pt, bottom=4pt, breakable]
\noindent\small\textbf{Задача №64.}\quad У прямокутному $\triangle CKD$: $\angle CDK=30^\circ$, катет $CK=2$ лежить навпроти $30^\circ$, тож $CD=\dfrac{CK}{\sin 30^\circ}=4$~(Б). $\angle KCD=60^\circ$; оскільки $CK$~--- бісектриса, $\angle ACK=60^\circ$, тоді у прямокутному $\triangle ACK$: $AK=CK\tan 60^\circ=2\sqrt{3}$, а $BC=AK=2\sqrt{3}$~(А). Відстань від $B$ до $AC$ дорівнює $\dfrac{2S_{ABC}}{AC}$; $AC=\dfrac{CK}{\cos 60^\circ}=4$, $S_{ABC}=\dfrac{1}{2}\cdot BC\cdot AB=\dfrac{1}{2}\cdot 2\sqrt{3}\cdot 2=2\sqrt{3}$, тож відстань $=\dfrac{2\cdot 2\sqrt{3}}{4}=\sqrt{3}$~(Г). Відповідь: \textbf{1--Б, 2--А, 3--Г}.
\par\vspace{0.1cm}\hfill\hyperlink{task:64}{\fbox{\footnotesize $\leftarrow$ Задача №64}}
\end{tcolorbox}
\par\vspace{0.3cm}\noindent\hypertarget{sol:67}{}
\begin{tcolorbox}[colback=titleBg!40, colframe=mainGreen!55, boxrule=0.4pt, arc=2pt, left=6pt, right=6pt, top=4pt, bottom=4pt, breakable]
\noindent\small\textbf{Задача №67.}\quad Нехай сторона квадрата $a$. Ламана $BADM=BA+AD+DM=a+a+\dfrac a2=\dfrac{5a}{2}=20$, звідки $a=8$ (\textbf{1}~--~Д). Висота трикутника до $BC$ дорівнює $\dfrac a2=4$, тож відстань від $K$ до $AD$: $a+\dfrac a2=12$ (\textbf{2}~--~А). У координатах $K(4;12)$, $M(8;4)$: $KM=\sqrt{4^2+8^2}=\sqrt{80}=4\sqrt5$ (\textbf{3}~--~Г). Відповідь: \textbf{1~--~Д, 2~--~А, 3~--~Г}.
\par\vspace{0.1cm}\hfill\hyperlink{task:67}{\fbox{\footnotesize $\leftarrow$ Задача №67}}
\end{tcolorbox}
\par\vspace{0.3cm}\noindent\hypertarget{sol:70}{}
\begin{tcolorbox}[colback=titleBg!40, colframe=mainGreen!55, boxrule=0.4pt, arc=2pt, left=6pt, right=6pt, top=4pt, bottom=4pt, breakable]
\noindent\small\textbf{Задача №70.}\quad Оскільки $BM\perp AD$ і $BC\parallel AD$, то $BM\perp BC$, тому в прямокутному трикутнику $BKC$: $BC=\sqrt{CK^2-BK^2}=\sqrt{17^2-8^2}=15$ см (1\,--\,Г). Далі $KM=\dfrac{3}{4}BK=6$, тож $BM=BK+KM=14$ см (2\,--\,В). Трикутники $NMK$ і $CBK$ подібні: $\dfrac{NM}{BC}=\dfrac{KM}{BK}=\dfrac{6}{8}$, звідки $NM=15\cdot\dfrac{6}{8}=11{,}25$ см (3\,--\,А). Відповідь: \textbf{1\,--\,Г, 2\,--\,В, 3\,--\,А}.
\par\vspace{0.1cm}\hfill\hyperlink{task:70}{\fbox{\footnotesize $\leftarrow$ Задача №70}}
\end{tcolorbox}
\par\vspace{0.3cm}\noindent\hypertarget{sol:73}{}
\begin{tcolorbox}[colback=titleBg!40, colframe=mainGreen!55, boxrule=0.4pt, arc=2pt, left=6pt, right=6pt, top=4pt, bottom=4pt, breakable]
\noindent\small\textbf{Задача №73.}\quad $KM=BM-BK=27-3=24$ см -- \textbf{Г}. Коло дотикається середин $AB$ і $CD$, тож його радіус дорівнює половині сторони $BC$; за симетрією $BC=BK+KM+BK=30$, отже $r=KO=15$ см -- \textbf{Б}. Відстань від $O$ до $BC$ дорівнює $\sqrt{r^2-(KM/2)^2}=\sqrt{225-144}=9$, тому $AB=2\cdot9=18$ см -- \textbf{В}. Відповідь: \textbf{1--Г, 2--Б, 3--В}.
\par\vspace{0.1cm}\hfill\hyperlink{task:73}{\fbox{\footnotesize $\leftarrow$ Задача №73}}
\end{tcolorbox}
\par\vspace{0.3cm}\noindent\hypertarget{sol:76}{}
\begin{tcolorbox}[colback=titleBg!40, colframe=mainGreen!55, boxrule=0.4pt, arc=2pt, left=6pt, right=6pt, top=4pt, bottom=4pt, breakable]
\noindent\small\textbf{Задача №76.}\quad $AM=AD+DM=12+2\cdot12=36$ см~--- \textbf{Г}. Нехай $H$~--- основа висоти рівнобедреного $\triangle DKM$ з вершини $K$; вона ділить $DM$ навпіл, тож $DH=12$, а $AH=AD+DH=24$. Оскільки $K$ лежить на прямій $AC$, прямокутні $\triangle ADC$ і $\triangle AHK$ подібні (спільний кут $A$, прямі кути при $D$ і $H$): $KH=CD\cdot\dfrac{AH}{AD}=16\cdot\dfrac{24}{12}=32$ см~--- \textbf{В}. З тієї самої подібності $AK=2\,AC$, де $AC=\sqrt{12^2+16^2}=20$, тож $CK=AK-AC=40-20=20$ см~--- \textbf{А}. Відповідь: \textbf{1~--~Г, 2~--~В, 3~--~А}.
\par\vspace{0.1cm}\hfill\hyperlink{task:76}{\fbox{\footnotesize $\leftarrow$ Задача №76}}
\end{tcolorbox}
\par\vspace{0.3cm}\noindent\hypertarget{sol:79}{}
\begin{tcolorbox}[colback=titleBg!40, colframe=mainGreen!55, boxrule=0.4pt, arc=2pt, left=6pt, right=6pt, top=4pt, bottom=4pt, breakable]
\noindent\small\textbf{Задача №79.}\quad $AC$~--- гіпотенуза рівнобедреного прямокутного $\triangle ABC$, тож катети $AB=BC=\dfrac{AC}{\sqrt2}=\dfrac{12}{\sqrt2}=6\sqrt2$~см (А). Оскільки $AN=NC=CK=\dfrac{AC}{2}=6$, маємо $AK=AN+NC+CK=18$~см (Г). Точка $B$ лежить на відстані катета від $C$; $BN$ обчислюється з прямокутного трикутника $BCN$ ($BC=6\sqrt2$, $CN=6$): тоді й $BM=\sqrt{BC^2+CM^2}$ дає $6\sqrt5$~см (Б). Відповідь: \textbf{1~--~Г, 2~--~А, 3~--~Б}.
\par\vspace{0.1cm}\hfill\hyperlink{task:79}{\fbox{\footnotesize $\leftarrow$ Задача №79}}
\end{tcolorbox}
\par\vspace{0.3cm}\noindent\hypertarget{sol:82}{}
\begin{tcolorbox}[colback=titleBg!40, colframe=mainGreen!55, boxrule=0.4pt, arc=2pt, left=6pt, right=6pt, top=4pt, bottom=4pt, breakable]
\noindent\small\textbf{Задача №82.}\quad Кути правильного трикутника по $60^\circ$, а в прямокутних рівнобедрених трикутниках $AKB$, $BMC$ кути при основі по $45^\circ$. Тоді $1$: $\angle ACM=\angle ACB+\angle BCM=60^\circ+45^\circ=105^\circ$ (В); $2$: $\angle KBM=\angle KBA+\angle ABC+\angle CBM=45^\circ+60^\circ+45^\circ=150^\circ$ (Д); $3$: $\angle KAM=75^\circ$ (А). Відповідь: \textbf{1~--~В, 2~--~Д, 3~--~А}.
\par\vspace{0.1cm}\hfill\hyperlink{task:82}{\fbox{\footnotesize $\leftarrow$ Задача №82}}
\end{tcolorbox}
\chapterTitle{ВІДПОВІДІ}
\noindent\small Відповіді згруповано за підтемами.\par\vspace{0.3cm}
\ansType{Завдання з вибором однієї відповіді (А--Д)}
\ansTheme{Кути, відрізки та обчислення в трикутниках}
\begin{multicols}{3}\noindent
\ansitem{В}
\ansitem{Б}
\ansitem{Д}
\ansitem{А}
\ansitem{Б}
\ansitem{Г}
\ansitem{А}
\ansitem{В}
\ansitem{В}
\ansitem{Г}
\ansitem{Г}
\ansitem{А}
\ansitem{Г}
\ansitem{В}
\ansitem{А}
\ansitem{Д}
\ansitem{Б}
\ansitem{Г}
\end{multicols}
\ansTheme{Чотирикутники: довжини й площі}
\begin{multicols}{3}\noindent
\ansitem{Д}
\ansitem{В}
\ansitem{Б}
\ansitem{Д}
\ansitem{Б}
\ansitem{Б}
\ansitem{В}
\ansitem{В}
\ansitem{А}
\ansitem{А}
\ansitem{Г}
\ansitem{В}
\ansitem{Д}
\ansitem{Б}
\ansitem{А}
\end{multicols}
\ansTheme{Площі складених фігур та кільця}
\begin{multicols}{3}\noindent
\ansitem{А}
\ansitem{Д}
\ansitem{А}
\ansitem{Г}
\ansitem{В}
\ansitem{Г}
\ansitem{Д}
\ansitem{Г}
\ansitem{А}
\end{multicols}
\ansTheme{Твердження та властивості геометричних фігур}
\begin{multicols}{3}\noindent
\ansitem{А}
\ansitem{Г}
\ansitem{Д}
\ansitem{В}
\ansitem{В}
\ansitem{Д}
\ansitem{Г}
\ansitem{В}
\ansitem{Б}
\ansitem{В}
\ansitem{А}
\ansitem{В}
\ansitem{Г}
\ansitem{Б}
\ansitem{Б}
\ansitem{Д}
\ansitem{Д}
\ansitem{Б}
\ansitem{В}
\ansitem{Б}
\ansitem{Г}
\end{multicols}
\ansType{Завдання на встановлення відповідності}
\ansTheme{Властивості геометричних фігур (відповідність)}
\begin{multicols}{3}\noindent
\ansitem{1~--~Б, 2~--~А, 3~--~Г}
\ansitem{1~--~Г, 2~--~А, 3~--~Б}
\ansitem{1~--~Б, 2~--~В, 3~--~Г}
\ansitem{1~--~Д, 2~--~А, 3~--~Г}
\ansitem{1~--~Д, 2~--~А, 3~--~Г}
\ansitem{1~--~А, 2~--~Г, 3~--~В}
\ansitem{1~--~Г, 2~--~В, 3~--~А}
\ansitem{1~--~Б, 2~--~В, 3~--~А}
\ansitem{1~--~А, 2~--~Б, 3~--~В}
\ansitem{1~--~Г, 2~--~Б, 3~--~В}
\ansitem{1~--~В, 2~--~Г, 3~--~А}
\ansitem{1~--~Г, 2~--~Д, 3~--~А}
\ansitem{1~--~Г, 2~--~В, 3~--~А}
\ansitem{1~--~В, 2~--~А, 3~--~Б}
\ansitem{1~--~Г, 2~--~Д, 3~--~В}
\ansitem{1~--~Г, 2~--~А, 3~--~Б}
\ansitem{1~--~Г, 2~--~А, 3~--~Б}
\ansitem{1~--~Г, 2~--~Б, 3~--~А}
\ansitem{1~--~В, 2~--~Д, 3~--~А}
\ansitem{1~--~А, 2~--~Б, 3~--~В}
\ansitem{1~--~Б, 2~--~А, 3~--~Д}
\end{multicols}
\end{document}
